\exercisesfor{5}

The conventions of § 5 remain valid unless otherwise mentioned.

\begin{enumerate}
\def\labelenumi{\arabic{enumi}.}
\item
  Let g be the 2-dimensional non-commutative solvable Lie algebra. Show that the Killing form of g is non-zero, that every invariant bilinear form on g is degenerate and that every derivation of g is inner.
\item
  \begin{enumerate}
  \def\labelenumii{(\alph{enumii})}
  \tightlist
  \item
    Show that in the 3-dimensional solvable Lie algebra over \(\mathbf{R}\) defined by the multiplication table \([x, y] = z\) , \([x, z] = -y\) , \([y, z] = 0\) , there exists no decreasing sequence of ideals of dimensions 3, 2, 1, 0.
  \end{enumerate}
\end{enumerate}

\begin{enumerate}
\def\labelenumi{(\alph{enumi})}
\setcounter{enumi}{1}
\tightlist
\item
  Show that in the non-commutative 2-dimensional solvable Lie algebra g there exists a sequence of ideals of dimensions 2, 1, 0 but that g is not nil-potent.
\end{enumerate}

\begin{enumerate}
\def\labelenumi{\arabic{enumi}.}
\setcounter{enumi}{2}
\item
  Let \(\mathfrak{g}\) be a solvable Lie algebra such that the conditions \(x\in \mathfrak{g},y\in \mathfrak{g}\) , \([[x,y],y] = 0\) imply \([x,y] = 0\) . Show that \(\mathfrak{g}\) is commutative. (Let \(k\) be the largest integer such that \(\mathcal{D}^{k - 1}\mathfrak{g}\neq \{0\}\) , \(\mathcal{D}^k\mathfrak{g} = \{0\}\) . Assuming \(k\geqslant 2\) , show first that \([\mathcal{D}^{k - 2}\mathfrak{g},\mathcal{D}^{k - 1}\mathfrak{g}] = \{0\}\) and then that \([\mathcal{D}^{k - 2}\mathfrak{g},\mathcal{D}^{k - 2}\mathfrak{g}] = \{0\}\) , whence a contradiction.)
\item
  Show that the centre of \(\mathfrak{st}(n,\mathbf{K})\) is zero and that that of \(\mathfrak{n}(n,\mathbf{K})\) is of dimension 1.
\item
  Let \(\mathfrak{g}\) be a Lie algebra and \((\mathcal{D}^0\mathfrak{g},\mathcal{D}^1\mathfrak{g},\ldots ,\mathcal{D}^n\mathfrak{g})\) the sequence of derived algebras of \(\mathfrak{g}\) ( \(n\geqslant 0,\mathcal{D}^{n - 1}\mathfrak{g}\neq \mathcal{D}^n\mathfrak{g}\) ). Then \(\dim \mathcal{D}^i\mathfrak{g} / \mathcal{D}^{i + 1}\mathfrak{g}\geqslant 2^{i - 1} + 1\) for \(1\leqslant i\leqslant n - 2\) . (Taking quotients by \(\mathcal{D}^n\mathfrak{g}\) , reduce it to the case where \(\mathfrak{g}\) is solvable. Then use the fact that \(\mathcal{D}_{\mathfrak{g}}\) is nilpotent and Exercise 7 (c) of § 4.)
\item
  \begin{enumerate}
  \def\labelenumii{(\alph{enumii})}
  \tightlist
  \item
    Verify that the following multiplication table defines a 5-dimensional solvable Lie algebra g:
  \end{enumerate}
\end{enumerate}

\[
\begin{array}{r} \left[ x _ {1}, x _ {2} \right] = x _ {5}, \qquad \left[ x _ {1}, x _ {3} \right] = x _ {3}, \qquad \left[ x _ {2}, x _ {4} \right] = x _ {4} \\ \left[ x _ {1}, x _ {4} \right] = \left[ x _ {2}, x _ {3} \right] = \left[ x _ {3}, x _ {4} \right] = \left[ x _ {5}, \mathfrak {g} \right] = 0. \end{array}
\]

\begin{enumerate}
\def\labelenumi{(\alph{enumi})}
\setcounter{enumi}{1}
\item
  Show that the orthogonal of \(\mathfrak{g}\) with respect to the Killing form is \(\mathcal{D}\mathfrak{g} = \mathrm{K}x_3 + \mathrm{K}x_4 + \mathrm{K}x_5\) . Deduce that \(\mathcal{D}\mathfrak{g}\) is the largest nilpotent ideal of \(\mathfrak{g}\) .
\item
  Show that there exists no supplementary subalgebra of \(\mathcal{D}\mathfrak{g}\) in \(\mathfrak{g}\) . Deduce that \(\mathfrak{g}\) is not the semi-direct product of a commutative algebra and a nilpotent ideal. (Show that this nilpotent ideal would necessarily be \(\mathcal{D}\mathfrak{g}\) .)
\end{enumerate}

\begin{enumerate}
\def\labelenumi{\arabic{enumi}.}
\setcounter{enumi}{6}
\item
  Let \(\mathfrak{g}\) be the 3-dimensional solvable Lie algebra with basis \((x,y,z)\) such that \([x,y] = z\) , \([x,z] = y\) , \([y,z] = 0\) . Show that the linear mapping which maps \(x\) to \(-x, y\) to \(-z, z\) to \(y\) is an automorphism of \(\mathfrak{g}\) of order 4. Compare this result with Exercise 21 (c) of §4.
\item
  \begin{enumerate}
  \def\labelenumii{(\alph{enumii})}
  \tightlist
  \item
    Let \(\mathfrak{g}_0\) be a 3-dimensional solvable real Lie algebra such that \(\mathcal{D}\mathfrak{g}_0\) is commutative and 2-dimensional. Let \(\mathfrak{g}\) be the algebra derived from \(\mathfrak{g}_0\) by extending the base field from \(\mathbf{R}\) to \(\mathbf{C}\) . For \(x \in \mathfrak{g}\) , let \(u_x\) be the restriction of \(\mathrm{ad}_{\mathfrak{g}}x\) to \(\mathcal{D}\mathfrak{g}\) . Show that the eigenvalues of \(u_x\) either have the same absolute value or are linearly dependent over \(\mathbf{R}\) . (We have \(x = \lambda y + z\) with \(z \in \mathcal{D}\mathfrak{g}\) , \(y \in \mathfrak{g}_0\) , \(\lambda \in \mathbf{C}\) and hence \(u_x = \lambda u_y\) . Now \(u_y\) is the \(\mathbf{C}\) -linear extension to \(\mathcal{D}\mathfrak{g}\) of an \(\mathbf{R}\) -linear endomorphism of \(\mathcal{D}\mathfrak{g}_{0}\) .)
  \end{enumerate}
\end{enumerate}

\begin{enumerate}
\def\labelenumi{(\alph{enumi})}
\setcounter{enumi}{1}
\item
  Show that there exists a 3-dimensional solvable complex Lie algebra g, with \(\mathcal{D}_{\mathfrak{g}}\) commutative and 2-dimensional, and an element \(x\) of g such that the restriction of \(\mathrm{ad}_{\mathfrak{g}}x\) to \(\mathcal{D}_{\mathfrak{g}}\) has eigenvalues which neither have the same absolute value nor are linearly dependent over R. (Construct g as a semi-direct product of a 1-dimensional algebra and a 2-dimensional commutative algebra.)
\item
  Show that the algebra constructed in (b) cannot be derived from a real Lie algebra by extending the base field from \(\mathbf{R}\) to \(\mathbf{C}\) .
\end{enumerate}

\begin{enumerate}
\def\labelenumi{\arabic{enumi}.}
\setcounter{enumi}{8}
\tightlist
\item
  Let \(\mathfrak{g}\) be a Lie algebra, \(\mathfrak{r}\) its radical, \(\mathfrak{n}\) its largest nilpotent ideal and D a derivation of \(\mathfrak{g}\) . Show that \(\mathcal{D}\mathfrak{g} \cap \mathfrak{r} \subset \mathfrak{n}\) . (Let \(\mathfrak{d}\) be the Lie algebra of derivations of \(\mathfrak{g}\) and \(\mathfrak{r}'\) its radical. If \(x \in \mathfrak{g}\) is such that \(\mathrm{D}x \in \mathfrak{r}\) , then
\end{enumerate}

\[
\operatorname{ad} (\mathrm{D} x) = [ \mathrm{D}, \operatorname{ad} x ] \in \mathscr {D} \mathfrak {d} \cap \mathfrak {r} ^ {\prime}
\]

(Corollary 2 to Proposition 5), hence \(\operatorname{ad}(\mathbf{D}x)\) is nilpotent (Theorem 1) and hence \(\mathbf{D}x \in \mathfrak{n}\) .)

\begin{enumerate}
\def\labelenumi{\arabic{enumi}.}
\setcounter{enumi}{9}
\item
  Let g be a Lie algebra, r its radical and a a subinvariant subalgebra of g. Show that the radical of a is \(a \cap r\) . (Apply Corollary 3 to Proposition 5 several times.)
\item
  Let \(\mathfrak{g}\) be a Lie algebra, \(\rho\) a finite-dimensional representation of \(\mathfrak{g}\) and \(\mathbb{E}\) the associative endomorphism algebra generated by 1 and \(\rho(\mathfrak{g})\) . Show that the largest nilpotency ideal \(n\) of \(\rho\) is equal to the set \(n'\) of \(x \in \mathfrak{g}\) such that \(\operatorname{Tr}(\rho(x)u) = 0\) for all \(u \in E\) . (To show that \(n' \subset n\) , show that \(n'\) is an ideal and that for all \(x \in n'\) the semi-simple component of \(\rho(x)\) is zero, by noting that \(\operatorname{Tr}((\rho(x))^n) = 0\) for every integer \(n > 0\) .)
\item
  Suppose that K is algebraically closed. Let g be a nilpotent Lie K-algebra. Let \(\rho\) be a finite-dimensional representation of g on a vector space V. For every linear form \(\lambda\) on g, let \(V^{\lambda}\) be the vector subspace of V consisting of the \(\xi \in V\) such that, for all \(x \in g\) , \((\rho(x) - \lambda(x)\mathrm{I})^{n}\xi = 0\) for sufficiently large n.
\end{enumerate}

\begin{enumerate}
\def\labelenumi{(\alph{enumi})}
\item
  The \(V^{\lambda}\) are stable under \(\rho(\mathfrak{g})\) and the sum of the \(V^{\lambda}\) is direct. (Use Exercise 22 of § 4.)
\item
  We have \(\mathrm{V} = \sum \mathrm{V}^{\lambda}\) . (If each \(\rho(x)\) has a single eigenvalue, \(\mathrm{V} = \mathrm{V}^{\lambda_0}\) by Corollary 2 to Theorem 1. If \(\rho(x_0)\) has at least two distinct eigenvalues, \(\mathrm{V}\) )
\end{enumerate}

is the direct sum of two non-trivial subspaces which are stable under \(\rho(\mathfrak{g})\) . Then argue by induction on the dimension of V.)

\begin{enumerate}
\def\labelenumi{\arabic{enumi}.}
\setcounter{enumi}{12}
\item
  Let \(\mathfrak{g}\) be a Lie Algebra and \(\mathfrak{D}\) the Lie algebra of derivations of \(\mathfrak{g}\) . For \(\mathfrak{g}\) to be characteristically nilpotent, it is necessary and sufficient that \(\mathfrak{D}\) be nilpotent and that \(\dim \mathfrak{g} > 1\) . (To see that the condition is sufficient, write \(\mathfrak{g}\) as a sum of subspaces \(\mathfrak{g}^{\lambda}\) , applying Exercise 12 to the identity mapping of \(\mathfrak{D}\) . Show that \([\mathfrak{g}^{\lambda}, \mathfrak{g}^{\mu}] \subset \mathfrak{g}^{\lambda + \mu}\) and that each \(\mathfrak{g}^{\lambda}\) is an ideal of \(\mathfrak{g}\) . Deduce that \(\mathfrak{g}^{\lambda}\) is commutative for \(\lambda \neq 0\) . Using again the fact that \(\mathfrak{D}\) is nilpotent, show that \(\mathfrak{g} = \mathfrak{g}^{0}\) if \(\dim \mathfrak{g} > 1\) . Otherwise \(\mathfrak{g} = \mathfrak{g}^{0} \times \mathfrak{h}\) , where \(\mathfrak{h}\) is commutative. Show first that \(\dim \mathfrak{h} \leqslant 1\) . If \(\dim \mathfrak{h} = 1\) , note that there would exist a derivation D of \(\mathfrak{g}\) such that D( \(\mathfrak{g}^{0}\) ) = \{0\} and D( \(\mathfrak{h}\) ) is contained in the centre of \(\mathfrak{g}^{0}\) ( \(\S 4\) , Exercise 8 ( \(a\) )).)
\item
  Let V be a finite-dimensional vector space over K and z an endomorphism of V. We adopt the notation \(z_{q}^{p}\) of Exercise 3 of §3. An endomorphism \(z'\) of V is called a replica of z if for all p and q every zero of \(z_{q}^{p}\) is a zero of \(z_{q}^{\prime p}\) . Show that if \(\operatorname{Tr}(zz') = 0\) for every replica \(z'\) of \(z\) , then z is nilpotent. (Use the proof of Lemma 3. In the notation of this proof, prove in particular that t is a replica of z.)
\item
  Let K be a field of characteristic 2. The identity representation of the nilpotent Lie algebra \(\mathfrak{sl}(2,\mathbf{K})\) on \(K^{2}\) defines a semi-direct product h of \(\mathfrak{sl}(2,\mathbf{K})\) by \(K^{2}\) . Show that h is solvable but that Dh is not nilpotent. Deduce that h admits no faithful linear representation by triangular matrices. Show also that the conclusions of Exercise 5 are false.
\item
  Let g be a Lie algebra over an arbitrary field K, A a commutative associative algebra over K and \(g' = g \otimes_{K} A\) , which can be considered as a Lie algebra over K.
\end{enumerate}

\begin{enumerate}
\def\labelenumi{(\alph{enumi})}
\item
  If D is a derivation of A, show that there exists one and only one derivation \(D'\) of \(g'\) such that \(\mathrm{D}'(x \otimes a) = x \otimes \mathrm{D}a\) for \(x \in g, a \in A\) .
\item
  Let \(p\) be a prime number, \(G\) a cyclic group of order \(p\) and \(s\) a generator of \(G\) . Suppose that \(K\) is of characteristic \(p\) and henceforth let \(A\) be the algebra of \(G\) over \(K\) . Show that there exists one and only one derivation \(D\) of \(A\) such that \(D(s^k) = ks^{k-1}\) for all \(k \in \mathbf{Z}\) . Show that the \(K\) -linear combinations of the \(x \otimes (s - 1)^k\) ( \(k = 1, 2, \ldots, p - 1, x \in \mathfrak{g}\) ) form a solvable ideal \(\tau\) of \(g'\) and that \(g'/\tau\) is isomorphic to \(\mathfrak{g}\) .
\item
  Take \(\mathfrak{g}\) to be simple (cf.~§6, no. 2, Definition 2). Then \(\mathfrak{r}\) is the radical of \(\mathfrak{g}'\) but is not a characteristic ideal. (Observe that \(\mathrm{D}(x\otimes (s - 1)) = 1\otimes x.\) )
\end{enumerate}

\begin{enumerate}
\def\labelenumi{\arabic{enumi}.}
\setcounter{enumi}{16}
\tightlist
\item
  Let g be a Lie algebra. Suppose that for every simple g-module M of finite dimension over K the \(x_{M}\) are permutable with one another. Show that g is solvable. (Observe that Dg is contained in the nilpotent radical and hence is solvable.)
\end{enumerate}

\exercisesfor{6}

The conventions of § 6 remain valid unless otherwise mentioned.

¶ 1. Let g be a semi-simple Lie algebra and ρ a finite-dimensional representation of g on M.

\begin{enumerate}
\def\labelenumi{(\alph{enumi})}
\item
  If \(\rho\) is simple and non-zero, \(\mathrm{H}^p (\mathfrak{g},\mathbf{M}) = \{0\}\) for all \(p\) . (Use Proposition 1 and Exercise 12 (j) of §3.)
\item
  For all \(\rho\) , \(\mathbf{H}^1(\mathfrak{g}, \mathbf{M}) = \{0\}\) . (If \(\rho\) is simple and non-zero, apply (a). If \(\rho\) is zero, use the fact that \(\mathfrak{g} = \mathcal{D}\mathfrak{g}\) . In the general case, argue by induction on the dimension of \(\rho\) : if N is a sub-g-module of M distinct from \(\{0\}\) and M, use the exact sequence:
\end{enumerate}

\[
\mathrm{H} ^ {1} (\mathfrak {g}, \mathrm{N}) \rightarrow \mathrm{H} ^ {1} (\mathfrak {g}, \mathrm{M}) \rightarrow \mathrm{H} ^ {1} (\mathfrak {g}, \mathrm{M} / \mathrm{N})
\]

established in Exercise 12 (f) of § 3.) Hence recover Remark 2 of no. 2.

\begin{enumerate}
\def\labelenumi{(\alph{enumi})}
\setcounter{enumi}{2}
\item
  Deduce from (b) a proof of Theorem 2. (Use Exercise 12 (g) of §3.) Deduce also from (b) that every derivation of g is inner. (Use Exercise 12 (h) of §3.)
\item
  For all \(\rho\) , \(\mathrm{H}^2(\mathfrak{g}, \mathrm{M}) = \{0\}\) . (Arguing as for (b), it suffices to consider the case where \(\rho = 0\) . Let \(c \in \mathrm{H}^2(\mathfrak{g}, \mathrm{M})\) . Consider, following Exercise 12 (i) of § 3, the central extension \(\mathfrak{h}\) of \(\mathfrak{g}\) by \(\mathrm{M}\) defined by \(c\) . The adjoint representation of \(\mathfrak{h}\) defines a representation of \(\mathfrak{g}\) on \(\mathfrak{h}\) . By Theorem 2, \(\mathrm{M}\) admits in \(\mathfrak{h}\) a supplement which is stable under \(\mathfrak{g}\) and hence the extension is trivial. Hence \(c = 0\) by Exercise 12 (i) of § 3.)
\item
  Deduce from (b) and (d) a proof of Theorem 5. (As in the text, it can be reduced to the case where the radical is commutative.)
\end{enumerate}

\begin{enumerate}
\def\labelenumi{\arabic{enumi}.}
\setcounter{enumi}{1}
\item
  Let g be a Lie algebra, r its radical and \((a_{0}, a_{1}, \ldots)\) a sequence of ideals of g defined as follows: (1) \(a_{0} = \{0\}\) ; (2) \(a_{i+1}/a_{i}\) is a maximal commutative ideal of \(g/a_{i}\) . Let p be the smallest integer such that \(a_{p} = a_{p+1} = \cdots\) . Show that \(r = a_{p}\) . (Show that \(g/a_{p}\) is semi-simple.)
\item
  For a Lie algebra g to be semi-simple, it is necessary and sufficient that it be reductive in every Lie algebra containing g as a subalgebra. (If g satisfies this condition, let M be a g-module of finite dimension over K. Considering M as a commutative algebra, form the semi-direct product of g and M, in which g is reductive. Deduce that M is semi-simple.)
\item
  Let g be a semi-simple Lie algebra and \(\rho\) a non-zero simple representation of g on a finite-dimensional space M. Let h be the corresponding semi-direct product. Show that \(h = Dh\) , that the centre of h is zero and that h is not the product of a semi-simple algebra and a solvable algebra.
\item
  Let \(\mathfrak{a}\) be a Lie algebra and \(\mathfrak{r}\) its radical. If \(\mathfrak{r}\) has a decreasing sequence of characteristic ideals \(\mathfrak{r} = \mathfrak{r}_0 \supset \mathfrak{r}_1 \supset \dots \supset \mathfrak{r}_n = \{0\}\) such that \(\dim \mathfrak{r}_i / \mathfrak{r}_{i+1} = 1\) for \(0 \leqslant i < n\) , \(\mathfrak{g}\) is the product of a semi-simple algebra and a solvable algebra.
\end{enumerate}

(Let s be a Levi subalgebra of g. For all \(x \in s\) , let \(\rho(x)\) be the restriction of \(ad_{g}x\) to r. Then \(\rho\) is a direct sum of 1-dimensional representations, which are zero because \(s = Ds\) .)

\begin{enumerate}
\def\labelenumi{\arabic{enumi}.}
\setcounter{enumi}{5}
\item
  Let \(\mathfrak{g}\) be a Lie algebra and \(\mathcal{D}^{\infty}\mathfrak{g}\) the intersection of the \(\mathcal{D}^{p}\mathfrak{g}\) ( \(p = 1,2,\ldots\) ). Show that \(\mathfrak{g}/\mathcal{D}^{\infty}\mathfrak{g}\) is solvable and that \(\mathcal{D}(\mathcal{D}^{\infty}\mathfrak{g}) = \mathcal{D}^{\infty}\mathfrak{g}\) . For \(\mathfrak{g}\) to be isomorphic to the product of a semi-simple algebra and a solvable algebra, it is necessary and sufficient that \(\mathcal{D}^{\infty}\mathfrak{g}\) be semi-simple.
\item
  \begin{enumerate}
  \def\labelenumii{(\alph{enumii})}
  \tightlist
  \item
    Let g be a Lie algebra, b a semi-simple ideal of g and a the centralizer of b in g, so that g is identified with \(b \times a\) . Show that, for every ideal \(\ell\) of g, then \(\ell = (\ell \cap \mathfrak{b}) \times (\ell \cap \mathfrak{a})\) . (Let \(\ell_{1}\) be the canonical projection of \(\ell\) onto b; it is an ideal of b and hence \(\mathcal{D}\ell_{1} = \ell_{1}\) ; deduce that \(\ell_{1} \subset \ell\) .)
  \end{enumerate}
\end{enumerate}

\begin{enumerate}
\def\labelenumi{(\alph{enumi})}
\setcounter{enumi}{1}
\item
  Let \(\beta\) be an invariant bilinear form on \(\mathfrak{g}\) . Show that \(\mathfrak{b}\) and \(\mathfrak{a}\) are orthogonal with respect to \(\beta\) (use the fact that \(\mathfrak{b} = \mathcal{D}\mathfrak{b}\) and that \(\beta = \beta_{1} + \beta_{2}\) , where \(\beta_{1}\) (resp. \(\beta_{2}\) ) is an invariant bilinear form whose restriction to \(\mathfrak{a}\) (resp. \(\mathfrak{b}\) ) is zero.
\item
  Deduce from (a) that there exists in \(\mathfrak{g}\) a largest semi-simple ideal. (Consider a maximal semi-simple ideal of \(\mathfrak{g}\) .)
\end{enumerate}

\begin{enumerate}
\def\labelenumi{\arabic{enumi}.}
\setcounter{enumi}{7}
\tightlist
\item
  Let g be a Lie algebra. An ideal h of g is called minimal if \(h \neq \{0\}\) and every ideal of g contained in h is equal to \(\{0\}\) or h.
\end{enumerate}

\begin{enumerate}
\def\labelenumi{(\alph{enumi})}
\item
  Every simple ideal of g is minimal.
\item
  Let \(\mathfrak{h}\) be a minimal ideal of \(\mathfrak{g}\) and \(\mathfrak{r}\) the radical of \(\mathfrak{g}\) . Then either \(\mathfrak{h} \subset \mathfrak{r}\) , in which case \(\mathfrak{h}\) is abelian, or \(\mathfrak{h} \cap \mathfrak{r} = \{0\}\) , in which case \(\mathfrak{h}\) is simple. (Use the fact that the derived ideal of a Lie algebra and the simple components of a semi-simple Lie algebra are characteristic ideals.)
\end{enumerate}

\begin{enumerate}
\def\labelenumi{\arabic{enumi}.}
\setcounter{enumi}{8}
\item
  For a Lie algebra g to be reductive, it is necessary and sufficient that its centre c be equal to its largest nilpotent ideal. (If the condition holds, let r be the radical of g; Dr is contained in the centre of r, hence r is nilpotent and hence r = c.)
\item
  Let \(\mathfrak{g}\) be a Lie algebra such that the conditions \(x \in \mathfrak{g}, y \in \mathfrak{g}, [[x, y], y] = 0\) imply \([x, y] = 0\) . Show that \(\mathfrak{g}\) is reductive. (Show that the radical \(\mathfrak{r}\) of \(\mathfrak{g}\) is commutative using Exercise 3 of §5. Then show that \([\mathfrak{g}, \mathfrak{r}] = 0\) .)
\end{enumerate}

¶ 11. (a) Let g be a Lie algebra, r its radical, s a Levi subalgebra of g and m an ideal of g containing r. There exists an ideal t of s, supplementary to m in g, such that \([m, t] \subset r\).

\begin{enumerate}
\def\labelenumi{(\alph{enumi})}
\setcounter{enumi}{1}
\tightlist
\item
  Let g be a Lie algebra and a subinvariant subalgebra of g. Then there exists a composition series \(g = g_{0} \supset g_{1} \supset \cdots \supset g_{k} = a\) such that \(g_{i}\) is the direct sum of \(g_{i+1}\) and a subalgebra \(b_{i}\) which is either 1-dimensional and contained in the radical \(r(g_{i})\) of \(g_{i}\) or simple and such that \([b_{i}, g_{i+1}] \subset r(g_{i+1})\) . (Reduce it to the case where a is an ideal of g. Let \(g' = g/a\) and \(r'\) be the radical of \(g'\) . The algebra \(g'/r'\) is the product of its simple ideals \(a_{1}, a_{2}, \ldots, a_{c}\) .
\end{enumerate}

Take a composition series of \(\mathfrak{g}'\) consisting of a composition series of \(\mathbf{r}'\) whose successive quotients are 1-dimensional and the inverse images of the ideals \(a_1, a_1 \times a_2, \ldots, a_1 \times a_2 \times \cdots \times a_c\) . Then take the inverse image in \(\mathfrak{g}\) of this composition series.)

\begin{enumerate}
\def\labelenumi{\arabic{enumi}.}
\setcounter{enumi}{11}
\tightlist
\item
  Let g be a Lie algebra, r its radical and D a derivation of g.
\end{enumerate}

\begin{enumerate}
\def\labelenumi{(\alph{enumi})}
\tightlist
\item
  If \(\mathrm{D}(\mathfrak{r}) = \{0\} \), D is inner. (Let c be the centralizer of r in g. The adjoint representation of g defines a representation \(x^{*} \mapsto \rho(x^{*})\) of \(g^{*} = g/r\) on c.~On the other hand, \([\mathrm{D}(\mathfrak{g}), \mathfrak{r}] \subset \mathrm{D}([\mathfrak{g}, \mathfrak{r}]) + [\mathfrak{g}, \mathrm{D}(\mathfrak{r})] = \{0\}\) , hence \(\mathrm{D}(\mathfrak{g}) \subset c\) and hence D defines a linear mapping \(D^{*}: g^{*} \to c\) . Show that
\end{enumerate}

\[
\mathrm{D} ^ {*} ([ x ^ {*}, y ^ {*} ]) = \rho (x ^ {*}) \mathrm{D} ^ {*} y ^ {*} - \rho (y ^ {*}) \mathrm{D} ^ {*} x ^ {*}
\]

for all \(x^{*}, y^{*}\) in \(\mathfrak{g}^{*}\) . Deduce that there exists \(a \in \mathfrak{c}\) such that \(D^{*}x^{*} = \rho(x^{*})a\) for all \(x^{*} \in \mathfrak{g}^{*}\) (cf.~no. 2, Remark 2), whence \(Dx = [x, a]\) for all \(x \in \mathfrak{g}\) .)

\begin{enumerate}
\def\labelenumi{(\alph{enumi})}
\setcounter{enumi}{1}
\tightlist
\item
  If D coincides on r with an inner derivation of g, D is inner. (Apply (a).)
\end{enumerate}

\begin{enumerate}
\def\labelenumi{\arabic{enumi}.}
\setcounter{enumi}{12}
\item
  Let g be a Lie algebra over an algebraically closed field, r its radical and \(\rho\) a finite-dimensional simple representation of g. Then \(\rho(z)\) is scalar for \(z \in r\) . (Reduce it to the case where g is reductive. In this case, r is the centre of g.)
\item
  Let X be a nilpotent matrix in \(\mathfrak{gl}(n, \mathbf{K})\) . Show that there exist in \(\mathfrak{gl}(n, \mathbf{K})\) two other matrices Y, H such that \([H, X] = X\) , \([H, Y] = -Y\) , \([X, Y] = H\) . (Argue by induction on n, using the Jordan form of X; thus reduce it to the case where \(X = E_{12} + E_{23} + \cdots + E_{n-1,n}\) ; then take Y to be a combination of the \(E_{k,k-1}\) and H a combination of the \(E_{jj'}\) .)
\end{enumerate}

¶ 15. (a) Let X, Y be two matrices of \(\mathfrak{gl}(n, \mathrm{K})\) . Show that, if \([X, Y] = X\) , X is nilpotent. (For every polynomial \(f \in K[T]\) , observe that

\[
[ f (X), Y ] = X f ^ {\prime} (X).)
\]

\begin{enumerate}
\def\labelenumi{(\alph{enumi})}
\setcounter{enumi}{1}
\item
  Let \(\mathfrak{g}\) be a Lie algebra and \(h\) , \(x\) two elements of \(\mathfrak{g}\) such that \([h, x] = x\) and there exists \(z \in \mathfrak{g}\) with \([x, z] = h\) . Show that there exists \(y \in \mathfrak{g}\) such that \([x, y] = h\) and \([h, y] = -y\) . (Observe first that \([z, h] + z\) belongs to the centralizer \(\mathfrak{n}\) of \(x\) in \(\mathfrak{g}\) . Then show that \(\mathfrak{n}\) is stable under ad \(h\) and that the restriction to \(\mathfrak{n}\) of I - ad \(h\) is bijective; to do this, note that ad \(x\) is nilpotent, using \((a)\) , and prove that, if \(\mathfrak{g}_{\mathfrak{l}}\) is the image of \(\mathfrak{g}\) under \((\mathrm{ad}x)^{\mathfrak{l}}\) , I - ad \(h\) gives on taking quotients a bijection of \((\mathfrak{n} \cap \mathfrak{g}_{\mathfrak{l}-1}) / (\mathfrak{n} \cap \mathfrak{g}_{\mathfrak{l}})\) ; for this use the relation \([\mathrm{ad}z, (\mathrm{ad}x)^{\mathfrak{k}}] = -k(\mathrm{ad}x)^{\mathfrak{k}-1}(\mathrm{ad}h + \frac{k-1}{2}I).)\)
\item
  Let \(\mathfrak{g}\) be a subalgebra of \(\mathfrak{gl}(n, \mathbf{K})\) ; suppose that, for every nilpotent matrix \(X \in \mathfrak{g}\) , there exist two other matrices \(H, Y\) in \(\mathfrak{g}\) such that \([H, X] = X\) , \([H, Y] = -Y\) , \([X, Y] = H\) . Let \(\mathfrak{h}\) be a subalgebra of \(\mathfrak{g}\) such that there exists
\end{enumerate}

a subspace m supplementary to h in g and such that \([h, m] \subset m\) . Show that h has the same property as g for all its nilpotent matrices (use (b)).

\begin{enumerate}
\def\labelenumi{\arabic{enumi}.}
\setcounter{enumi}{15}
\item
  Let g be a subalgebra of \(\mathfrak{gl}(n,\mathbf{K})\) such that the identity representation of g is semi-simple. Show that every nilpotent matrix \(X\in g\) is contained in a 3-dimensional simple Lie subalgebra of g. (Show that g is reductive in \(\mathfrak{gl}(n,\mathbf{K})\) ; then use Exercises 14 and 15 (c).)
\item
  Show that the algebra derived from a real simple Lie algebra by extending the base field from \(\mathbf{R}\) to \(\mathbf{C}\) is not always simple. (Let \(\mathfrak{g}\) be a real simple Lie algebra. If \(\mathfrak{g}' = \mathfrak{g}_{(\mathbf{C})}\) is simple, let \(\mathfrak{h}\) be the real Lie algebra derived from \(\mathfrak{g}'\) by restricting the field of scalars to \(\mathbf{R}\) . We know that \(\mathfrak{h}\) is simple. Then \(\mathfrak{h}_{(\mathbf{C})}\) is, by Exercise 4 of §1, the product of two algebras isomorphic to \(\mathfrak{g}'\) . Cf. also Exercise 26 (b).
\item
  \begin{enumerate}
  \def\labelenumii{(\alph{enumii})}
  \tightlist
  \item
    Let g be a simple Lie algebra. Every invariant bilinear form on g is either zero or non-degenerate. If K is algebraically closed, every invariant bilinear form \(\beta\) on g is proportional to the Killing form \(\beta_{0}\) . (Consider the endomorphism \(\sigma\) of the vector space g defined by \(\beta(x,y)=\beta_{0}(\sigma(x),y)\) and show that \(\sigma\) commutes with \(ad_{g}x\) and is hence scalar.) Show that this result may be false if K is not algebraically closed. (Use Exercise 17 and the fact that the dimension of the space of invariant bilinear forms does not change when the base field is extended.)
  \end{enumerate}
\end{enumerate}

\begin{enumerate}
\def\labelenumi{(\alph{enumi})}
\setcounter{enumi}{1}
\item
  Let g be a semi-simple Lie algebra over an algebraically closed field. Deduce from (a) and Exercise 7 (b) that the dimension of the space of invariant bilinear forms on g is equal to the number of simple components of g and that all these forms are symmetric.
\item
  Let g be a simple Lie algebra, M the dual space of the vector space g, with the dual representation of the adjoint representation, and b the semidirect product of g and M defined by \(x \mapsto x_{M}\) . For y, z in g and \(y'\) , \(z'\) in M, we write \(\beta(y + y', z + z') = \langle y, z' \rangle + \langle z, y' \rangle\) . Show that \(\beta\) is a non-degenerate invariant symmetric bilinear form on b, which is associated with no representation of b. (Observe that the radical and nilpotent radical of b are equal to M.)
\end{enumerate}

\begin{enumerate}
\def\labelenumi{\arabic{enumi}.}
\setcounter{enumi}{18}
\tightlist
\item
  Let \(g\) be a reductive Lie algebra, \(U\) its enveloping algebra, \(Z\) the centre of \(U\) and \(V\) the subspace of \(U\) generated by the elements of the form \(uv - vu\) ( \(u \in U, v \in U\) ).
\end{enumerate}

\begin{enumerate}
\def\labelenumi{(\alph{enumi})}
\item
  U is the direct sum of Z and V. (Apply Proposition 6 of §3 to the representation \(x \mapsto \mathrm{ad}_U x\) of \(\mathfrak{g}\) on U, observing that, in the decomposition \(\mathrm{U} = \sum_{n} \mathrm{U}^n\) of §2, no. 7, Corollary 4 to Theorem 1, the \(\mathrm{U}^n\) are stable under this representation.)
\item
  Let \(u \mapsto u^{\natural}\) be the projector of U onto Z parallel to V. Show that \((uv)^{\natural} = (vu)^{\natural}, (zu)^{\natural} = zu^{\natural}\) for all \(u \in \mathrm{U}, v \in \mathrm{U}, z \in \mathrm{Z}\) .
\item
  Let R be a two-sided ideal of U. Show that \(\mathbf{R} = (\mathbf{R} \cap \mathbf{Z}) + (\mathbf{R} \cap \mathbf{V})\) . (To show that the component in V of an element r of R belongs to R, reduce it to the case where K is algebraically closed; observe that R is stable under the representation \(\rho\) of U on U which extends \(x \mapsto ad_{U}x\) ; decompose r into its components in the \(U^{n}\) ; finally, apply Corollary 2 to Theorem 1 of Algebra, Chapter VIII, § 4, no. 2 to the restriction of \(\rho\) to \(U^{n}\) .)
\item
  Let \(\mathbf{R}'\) be an ideal of \(Z\) . Let \(\mathbf{R}_1\) be the two-sided ideal of \(U\) generated by \(\mathbf{R}'\) . Show that \(\mathbf{R}_1 \cap Z = \mathbf{R}'\) . (Use (b).)
\end{enumerate}

\begin{enumerate}
\def\labelenumi{\arabic{enumi}.}
\setcounter{enumi}{19}
\tightlist
\item
  Suppose that \(\mathbf{K}\) is algebraically closed. Let \(\mathfrak{g}\) be a Lie algebra and \(\mathfrak{c}\) its centre.
\end{enumerate}

\begin{enumerate}
\def\labelenumi{(\alph{enumi})}
\item
  If \(\mathfrak{g}\) admits a finite-dimensional faithful simple representation, \(\mathfrak{g}\) is reductive and \(\dim \mathfrak{c} \leqslant 1\) .
\item
  Conversely, if g is reductive and dim \(c \leqslant 1\) , g admits a finite-dimensional faithful simple representation. (If g is simple, use the adjoint representation. If \(g = g_{1} \times g_{2}\) and \(g_{1}, g_{2}\) admit finite-dimensional faithful simple representations \(\rho_{1}, \rho_{2}\) on spaces \(M_{1}, M_{2}\) , show that
\end{enumerate}

\[
(x _ {1}, x _ {2}) \mapsto \rho_ {1} (x _ {1}) \otimes 1 + 1 \otimes \rho_ {2} (x _ {2})
\]

is a semi-simple representation \(\rho\) of \(g\) , then that the centralizer of the \(g\) -module \(\mathrm{M}_1 \otimes \mathrm{M}_2\) reduces to the scalars and hence that \(\rho\) is simple. When \(g_1\) and \(g_2\) are semi-simple or \(g_1\) is semi-simple and \(g_2\) is commutative, show that \(\rho\) is faithful by considering the kernel of \(\rho\) which is an ideal of \(g_1 \times g_2\) .

\begin{enumerate}
\def\labelenumi{\arabic{enumi}.}
\setcounter{enumi}{20}
\item
  Let V be a vector space over K of finite dimension n \textgreater{} 1. Show that \(\mathfrak{sl}(V)\) is simple. (Reduce it to the case where K is algebraically closed. Suppose that \(\mathfrak{sl}(V) = a \times b\) , with dim a = a \textgreater{} 0, dim b = b \textgreater{} 0. Let A (resp. B) be the associative algebra generated by 1 and a (resp. b). Then V can be considered as a simple \((\mathrm{A} \otimes_{\mathrm{K}} \mathrm{B})\) -module. Hence there exists an A-module P of finite dimension p over K and a B-module Q of finite dimension q over K such that V is \((\mathrm{A} \otimes_{\mathrm{K}} \mathrm{B})\) -isomorphic to \(P \otimes_{K} Q\) (Algebra, Chapter VIII, §7, no. 7, Proposition 8 and no. 4, Theorem 2); P and Q are faithful and hence \(a \leqslant p^{2}\) , \(b \leqslant q^{2}\) . On the other hand, \(a + b = n^{2} - 1\) and pq = n, whence \((p^{2} - 1)(q^{2} - 1) \leqslant 2\) , which is contradictory.)
\item
  \begin{enumerate}
  \def\labelenumii{(\alph{enumii})}
  \tightlist
  \item
    Let \(\mathfrak{g}\) be a nilpotent Lie algebra over an algebraically closed field \(\mathbf{K}\) of arbitrary characteristic. Let \(\mathfrak{z}\) be its centre, \(\rho\) a finite-dimensional representation of \(\mathfrak{g}\) and \(\beta\) the associated bilinear form. Show that \(\mathfrak{z}\cap \mathcal{D}\mathfrak{g}\) is orthogonal to \(\mathfrak{g}\) with respect to \(\beta\) . (Reduce it, using a Jordan-Hölder series, to the case where \(\rho\) is simple. Then note that, for \(z\in \mathfrak{g}\cap \mathcal{D}\mathfrak{g},\rho (z)\) is a scalar matrix with trace zero. Then use Exercise 22 (b) of §4.) Deduce that, if \(\beta\) is nondegenerate, \(\mathfrak{g}\) is commutative. (Use Exercise 3 (a) of §4.)
  \end{enumerate}
\end{enumerate}

\begin{enumerate}
\def\labelenumi{(\alph{enumi})}
\setcounter{enumi}{1}
\item
  Suppose that \(\mathbf{K}\) is of characteristic 2. Show that for the solvable Lie algebra \(\mathfrak{gl}(2,\mathbf{K}) = \mathfrak{g}\) the bilinear form associated with the identity representation is non-degenerate, but that the centre \(\mathfrak{z}\) is contained in \(\mathcal{D}\mathfrak{g}\) and is \(\neq \{0\}\) .
\item
  Let g be the 6-dimensional Lie algebra over K with basis \((a, b, c, d, e, f)\) and multiplication table \([a, b] = -[b, a] = d\) , \([a, c] = -[c, a] = e\) , \([b, c] = -[c, b] = f\) and the other brackets are 0. Let β be the bilinear form on g such that β \((c, d) = \beta(d, c) = 1\) , β \((a, f) = \beta(f, a) = 1\) , β \((b, e) = \beta(e, b) = -1\) and the other values of β at ordered pairs of the basis of g in question are 0. Show that β is invariant, g nilpotent, \(\mathfrak{z} = \mathcal{D}g \neq \{0\}\) and β non-degenerate.
\end{enumerate}

\begin{enumerate}
\def\labelenumi{\arabic{enumi}.}
\setcounter{enumi}{22}
\tightlist
\item
  Let \(\mathfrak{g}\) be a 3-dimensional non-commutative Lie algebra over a field \(\mathbf{K}\) of arbitrary characteristic.
\end{enumerate}

\begin{enumerate}
\def\labelenumi{(\alph{enumi})}
\item
  If \(\mathfrak{g}\) has centre \(\mathfrak{z} \neq \{0\}\) , then dim \(\mathfrak{z} = 1\) (§ 1, Exercise 2) and dim \(\mathcal{D}_{\mathfrak{g}} = 1\) . If \(\mathfrak{z} \neq \mathcal{D}_{\mathfrak{g}}, \mathfrak{g}\) is the product of \(\mathfrak{z}\) and the 2-dimensional non-commutative algebra. If \(\mathfrak{z} = \mathcal{D}_{\mathfrak{g}}, \mathfrak{g}\) is the 3-dimensional non-commutative nilpotent algebra (§ 4, Exercise 9 (a)).
\item
  If \(\mathfrak{z} = \{0\}\) but there exists in \(\mathfrak{g}\) a 2-dimensional commutative subalgebra, this subalgebra is unique and is equal to \(\mathcal{D}_{\mathfrak{g}}\) ; for all \(x \notin \mathcal{D}_{\mathfrak{g}}\) , the restriction \(u\) of \(ad x\) to \(\mathcal{D}_{\mathfrak{g}}\) is a bijection of this vector space, determined to within a multiplicative constant. Conversely, every automorphism \(u\) of a 2-dimensional vector space \(\mathrm{Ka} + \mathrm{Kb}\) determines a Lie algebra structure on \(\mathfrak{g} = \mathrm{Ka} + \mathrm{Kb} + \mathrm{Kc}\) by the conditions \([a, b] = 0\) , \([c, a] = u(a)\) , \([c, b] = u(b)\) . For two Lie algebras thus defined by automorphisms \(u_1, u_2\) to be isomorphic, it is necessary and sufficient that the matrices of \(u_1\) and \(u_2\) be similar to within a scalar factor.
\item
  Suppose that there exists in g no commutative subalgebra of dimension \textgreater1. Show that there then exists \(a \in g\) such that \(a \notin (\operatorname{ad} a)g\) (assuming that an element x does not have this property, show, using the Jacobi identity, that another element has the desired property). There then exists a basis \((a, b, c)\) of g such that \([a, b] = c\) , \([a, c] = \beta b\) , \([b, c] = \gamma a\) , with \(\beta \neq 0\) and \(\gamma \neq 0\) , and g is simple. Let \(\phi\) be a canonical isomorphism of g onto the exterior power \(\bigwedge^{2}\mathfrak{g}^{*}\) of the dual space of \(\mathfrak{g}\) , determined to within an invertible constant factor (Algebra, Chapter III, § 8, no. 5); let \(u\) be the linear mapping of \(\mathfrak{g}\) into \(\mathfrak{g}^{*}\) defined by the condition \(\langle [x,y],u(z)\rangle = \langle x\wedge y,\phi (z)\rangle\) ; the bilinear form \(\Phi (x,y) = \langle x,u(y)\rangle\) on \(\mathfrak{g}\times \mathfrak{g}\) is symmetric and non-degenerate and \(2\Phi\) is the Killing form of \(\mathfrak{g}\) to within a factor \(\neq 0\) . For two 3-dimensional simple Lie algebras over K to be isomorphic, it is necessary and sufficient that the corresponding bilinear forms be equivalent to within a constant factor \(\neq 0\) .
\item
  If \(\mathbf{K}\) is not of characteristic 2, the simple algebra \(\mathfrak{g}\) defined in (c) is isomorphic to the Lie algebra \(\mathfrak{o}(\Phi)\) of the formal orthogonal group \(\mathbf{G}(\Phi)\) (§ 1, Exercise 26 (a)). For the existence in \(\mathfrak{g}\) of vectors \(x\) such that ad \(x\) admits eigenvectors not proportional to \(x\) , it is necessary and sufficient that \(\Phi\) be of index \(>0\) ; \(\mathfrak{g}\) then admits a basis \((a, b, c)\) such that \([a, b] = b\) , \([a, c] = -c\) , \([b, c] = a\) .
\item
  If K is of characteristic 2, show that there is no 2-mapping on g and therefore that g is not the Lie algebra of a formal group. Show that g admits derivations which are not inner.
\end{enumerate}

¶ 24. Let K be a field of arbitrary characteristic p.~We again adopt Definition 1.

\begin{enumerate}
\def\labelenumi{(\alph{enumi})}
\tightlist
\item
  Show that, unless \(p = n = 2\) , the only non-trivial ideals of \(\mathfrak{gl}(n, \mathbf{K})\)
\end{enumerate}

are the 1-dimensional centre \(\mathfrak{z}\) with basis \(\sum_{i=1}^{n} E_{ii}\) and the Lie algebra \(\mathfrak{sl}(n, \mathrm{K})\) consisting of the matrices of trace 0. (Except in the exceptional case indicated, note that if an ideal \(\mathfrak{a}\) contains one of the \(E_{ij}\) , it necessarily contains \(\mathfrak{sl}(n, \mathrm{K})\) . If \(\mathfrak{a}\) contains an element not belonging to \(\mathfrak{b}\) , by multiplying this element by at most four suitably chosen elements \(E_{ij}\) , a non-zero multiple of one of the \(E_{ij}\) is obtained.) If \(n\) is not a multiple of \(p\) , show that \(\mathfrak{gl}(n, \mathrm{K})\) is the direct sum of \(\mathfrak{z}\) and \(\mathfrak{sl}(n, \mathrm{K})\) and \(\mathfrak{sl}(n, \mathrm{K})\) is simple. If on the contrary \(n\) is a multiple of \(p\) and \(n > 2\) , show that \(\mathfrak{sl}(n, \mathrm{K}) / \mathfrak{z}\) is simple (same methods).

\begin{enumerate}
\def\labelenumi{(\alph{enumi})}
\setcounter{enumi}{1}
\tightlist
\item
  Show that, for \(n\) a multiple of \(p\) and \(n > 2\) , \(\mathfrak{gl}(n, \mathbf{K}) / \mathfrak{sl}(n, \mathbf{K})\) has radical \(\{0\}\) , but admits the quotient \(\mathfrak{gl}(n, \mathbf{K}) / \mathfrak{sl}(n, \mathbf{K})\) which is Abelian.
\end{enumerate}

¶ 25. Let K be a field of arbitrary characteristic p.~Let \(\mathfrak{sp}(2n, \mathrm{K})\) denote the Lie algebra of the formal symplectic group in 2n variables over K (§ 1, Exercise 26 (c)). Show that this Lie algebra, which is identified with a sub-algebra of \(\mathfrak{gl}(2n, \mathrm{K})\) , has basis consisting of the elements

\[
\begin{array}{c} H _ {i} = E _ {i i} - E _ {i + n, i + n} \quad (1 \leqslant i \leqslant n), \\ F _ {i j} = E _ {i j} - E _ {j + n, i + n} \quad (1 \leqslant i, j \leqslant n, i \neq j), \\ G _ {i j} ^ {\prime} = E _ {i, j + n} + E _ {j, i + n}, \qquad G _ {i j} ^ {\prime \prime} = E _ {i + n, j} + E _ {j + n, i} \quad (1 \leqslant i <   j \leqslant n), \\ E _ {i, i + n} \quad \text { and } \quad E _ {i + n, i} \quad (1 \leqslant i \leqslant n). \end{array}
\]

\begin{enumerate}
\def\labelenumi{(\alph{enumi})}
\item
  Show that, if \(p \neq 2\) and \(n \geqslant 1\) , the algebra \(\mathfrak{sp}(2n, \mathrm{K})\) is simple (for \(n = 1\) , \(\mathfrak{sp}(2, \mathrm{K}) = \mathfrak{sl}(2, \mathrm{K})\) ). (Same method as in Exercise 24.)
\item
  If \(p = 2\) and \(n \geqslant 3\) , show that the \(H_{i}, F_{ij}, G_{ij}'\) and \(G_{ij}''\) form an ideal \(\mathfrak{a}\) of dimension \(n(2n - 1)\) ; the elements of \(\mathfrak{a}\) of the form
\end{enumerate}

\[
\sum_ {i = 1} ^ {n} \lambda_ {i} H _ {i} + \sum_ {i \neq j} \alpha_ {i j} F _ {i j} + \sum_ {i <   j} \beta_ {i j} G _ {i j} ^ {\prime} + \sum_ {i <   j} \gamma_ {i j} G _ {i j} ^ {\prime \prime}
\]

such that \(\sum_{i=1}^{n} \lambda_i = 0\) form an ideal \(b \subset a\) of dimension \(n(2n - 1) - 1\) and the multiples of \(\sum_{i=1}^{n} H_i\) an ideal \(c\) , the centre of \(\mathfrak{sp}(2n, K)\) ; \(b, c\) are the only ideals of \(\mathfrak{a},\mathfrak{b} = \mathcal{D}(\mathfrak{sp}(2n,\mathbf{K}))\) and \(\mathfrak{b} / \mathfrak{c}\) is simple. (Same method.) What are the ideals of \(\mathfrak{sp}(4,\mathbf{K})\) when \(\mathbf{K}\) is of characteristic 2?

¶ 26. Let K be a field of characteristic ≠2 and Φ a non-degenerate symmetric bilinear form on an n-dimensional vector space over K.

\begin{enumerate}
\def\labelenumi{(\alph{enumi})}
\tightlist
\item
  Show that, if \(n \geqslant 5\) , the Lie algebra \(\mathfrak{o}(\Phi)\) is simple. (By extending the base field, it can be reduced to the case where the index of \(\Phi\) is \(m = [n/2]\) . If for example \(n = 2m\) is even, \(\mathfrak{o}(\Phi)\) has basis consisting of the elements
\end{enumerate}

\[
\begin{array}{c} H _ {i} = E _ {i i} - E _ {i + m, i + m}, \quad F _ {i j} = E _ {i j} - E _ {j + m, i + m} \quad (1 \leqslant i, j \leqslant m, i \neq j), \\ G _ {i j} ^ {\prime} = E _ {i, j + m} - E _ {j, i + m} \quad \text { and } \quad G _ {i j} ^ {\prime \prime} = E _ {i + m, j} - E _ {j + m, i} \quad (1 \leqslant i <   j \leqslant m); \end{array}
\]

then argue as in Exercises 24 and 25. Proceed similarly when \(n = 2m + 1\) is odd.)

\begin{enumerate}
\def\labelenumi{(\alph{enumi})}
\setcounter{enumi}{1}
\tightlist
\item
  Let \(\Delta\) be the discriminant of \(\Phi\) with respect to a basis. Suppose that \(n = 4\) . Show that, if \(\Delta\) is not a square in \(\mathbf{K}\) , \(\mathfrak{o}(\Phi)\) is simple. If \(\Delta\) is a square, \(\mathfrak{o}(\Phi)\) is the product of two isomorphic 3-dimensional simple Lie algebras. (Use the structure of \(\mathbf{G}(\Phi)\) described in Algebra, Chapter IX, §9, Exercise 16.) Deduce an example of a simple Lie algebra which becomes non-simple when extending the base field.
\end{enumerate}

\begin{enumerate}
\def\labelenumi{\arabic{enumi}.}
\setcounter{enumi}{26}
\tightlist
\item
  \begin{enumerate}
  \def\labelenumii{(\alph{enumii})}
  \tightlist
  \item
    Show that, in a finite-dimensional Lie \(p\) -algebra, the radical is a \(p\) -ideal. (Use Exercise 22 (c) of § 1.)
  \end{enumerate}
\end{enumerate}

\begin{enumerate}
\def\labelenumi{(\alph{enumi})}
\setcounter{enumi}{1}
\tightlist
\item
  For a Lie \(p\) -algebra \(\mathfrak{g}\) to have zero radical, it is necessary and sufficient that \(\mathfrak{g}\) contain no commutative \(p\) -ideal \(\neq \{0\}\) . (Use Exercise 22 (c) of § 1.)
\end{enumerate}

\exercisesfor{7}

The conventions of § 7 remain valid unless otherwise mentioned.

\begin{enumerate}
\def\labelenumi{\arabic{enumi}.}
\tightlist
\item
  \begin{enumerate}
  \def\labelenumii{(\alph{enumii})}
  \tightlist
  \item
    For a Lie K-algebra to be nilpotent, it is necessary and sufficient that it be isomorphic to a subalgebra of an algebra \(n(n, \mathbf{K})\) .
  \end{enumerate}
\end{enumerate}

\begin{enumerate}
\def\labelenumi{(\alph{enumi})}
\setcounter{enumi}{1}
\tightlist
\item
  Suppose that K is algebraically closed. For a Lie K-algebra to be solvable, it is necessary and sufficient that it be isomorphic to a subalgebra of an algebra \(t(n, \mathbf{K})\) .
\end{enumerate}

\begin{enumerate}
\def\labelenumi{\arabic{enumi}.}
\setcounter{enumi}{1}
\item
  Let g be a Lie algebra and n its largest nilpotent ideal. There exist a finite-dimensional vector space V and an isomorphism of g onto a subalgebra of \(\mathfrak{sl}(V)\) , which maps every element of n to a nilpotent endomorphism of V. (Use Ado's Theorem and Exercise 5 of § 1.)
\item
  Let g be a Lie algebra and U its enveloping algebra.
\end{enumerate}

\begin{enumerate}
\def\labelenumi{(\alph{enumi})}
\item
  Show that for all \(u \in \mathrm{U}\) ( \(u \neq 0\) ) there exists a finite-dimensional representation \(\pi\) of \(\mathrm{U}\) such that \(\pi(u) \neq 0\) . (Use Exercise 2 and Exercise 5 (b) of § 3.)
\item
  Deduce from (a) that if g is semi-simple there exists an infinity of inequivalent finite-dimensional simple representations of g. (Let \(\rho_{1},\ldots,\rho_{n}\) be finite-dimensional simple representations of g, \(N_{1},\ldots,N_{n}\) the annihilators of the corresponding U-modules and \(N=\bigcap_{i=1}^{n}N_{i}\) . Show that N is of finite co-dimension in U. Let \(n\in N,n\neq0\) . Apply (a) to n.)
\end{enumerate}

\begin{enumerate}
\def\labelenumi{\arabic{enumi}.}
\setcounter{enumi}{3}
\tightlist
\item
  Let g be a Lie algebra, a a subinvariant subalgebra of g, \(\mathfrak{r}(a)\) the radical of a and \(\rho\) a finite-dimensional representation of a. For \(\rho\) to admit a finite-dimensional extension to g, it is necessary and sufficient that \(\mathcal{D}_{\mathfrak{g}} \cap \mathfrak{r}(a)\) be contained in the largest nilpotency ideal of \(\rho\) . (For the sufficiency proceed as follows; let \(g = g_{0} \supset g_{1} \supset \cdots \supset g_{k} = a\) be a decomposition series of g with the properties of Exercise 11 (b) of §6. Show by induction the existence of an extension \(\rho_{i}\) of \(\rho\) to \(g_{0}\) , whose largest nilpotency ideal contains \(\mathcal{D}_{\mathfrak{g}} \cap \mathfrak{r}(g_{i})\) . In the notation of Exercise 11 (b) of §6, the passage from \(\rho_{i+1}\) to \(\rho_{i}\) follows from Theorem 1 when \(b_{i}\) is simple, or when \(b_{i}\) is 1-dimensional and
\end{enumerate}

\[
\mathscr {D} \mathfrak {g} \cap \mathfrak {r} (\mathfrak {g} _ {i + 1}) = \mathscr {D} \mathfrak {g} \cap \mathfrak {r} (\mathfrak {g} _ {i}).
\]

When \(\mathcal{D}\mathfrak{g} \cap \mathfrak{r}(\mathfrak{g}_{i+1}) \neq \mathcal{D}\mathfrak{g} \cap \mathfrak{r}(\mathfrak{g}_i)\) , it can be assumed that \(\mathfrak{h}_i \subset \mathcal{D}\mathfrak{g} \cap \mathfrak{r}(\mathfrak{g}_i)\) . Then \(\mathfrak{r}(\mathfrak{g}_i) = \mathfrak{r}(\mathfrak{g}) \cap \mathfrak{g}_i\) (§ 5, Exercise 10) and Theorem 1 can be applied.)

¶ 5. Let K be a field of characteristic p \textgreater{} 0, g a Lie algebra of finite dimension n over K and U the enveloping algebra of g.

\begin{enumerate}
\def\labelenumi{(\alph{enumi})}
\item
  A \(p\) -polynomial (in one indeterminate) over \(\mathbf{K}\) is any polynomial in \(\mathrm{K}[\mathrm{X}]\) whose only terms \(\neq 0\) are of degree a power of \(p\) . Show that for every polynomial \(f(\mathbf{X}) \neq 0\) in \(\mathrm{K}[\mathrm{X}]\) there exists \(g(\mathbf{X}) \in \mathrm{K}[\mathbf{X}]\) such that \(f(\mathbf{X})g(\mathbf{X})\) is a \(p\) -polynomial (consider the remainders in Euclidean division of the \(\mathbf{X}^{p_k}\) by \(f(\mathbf{X})\) ).
\item
  Show that for every element \(z \in \mathfrak{g}\) there exists a non-zero polynomial \(f(\mathbf{X}) \in \mathbf{K}[\mathbf{X}]\) such that \(f(z)\) belongs to the centre C of U. (Consider the minimal polynomial of the endomorphism ad \(z\) and apply (a) to this polynomial, and also formula (1) of Exercise 19 of § 1.)
\item
  Let \((e_i)_{1\leqslant i\leqslant n}\) be a basis of \(\mathfrak{g}\) . For each \(i\) let \(f_{i}\neq 0\) be a polynomial in \(\mathrm{K}[\mathrm{X}]\) such that \(f_{i}(e_{i})\in \mathrm{C}\) ; let \(d_{i}\) be its degree and \(\mathrm{L}\) the two-sided ideal of \(\mathrm{U}\) generated by the \(y_{i} = f_{i}(e_{i})\) . Show that the classes mod. \(\mathrm{L}\) of the elements \(e_1^{\alpha_1}\ldots e_n^{\alpha_n}\) , where \(0\leqslant \alpha_{i} < d_{i}\) form a basis of \(\mathrm{U / L}\) .
\item
  Show that \(\mathfrak{g}\) admits a finite-dimensional faithful linear representation (choose the \(f_{i}\) so that \(d_{i} > 1\) for all \(i\) ).
\item
  Show that if \(g \neq \{0\}\) there exists a non-semi-simple finite-dimensional linear representation of g. (Assuming that the \(f_{t}\) are of degrees \(d_{t} > 1\) , replace the \(f_{t}\) by \(f_{t}^{2}\) in the definition of L and note that there then exist nilpotent elements \(\neq 0\) in the centre of U/L and hence that U/L is not semi-simple.)
\end{enumerate}


\chapter{Free Lie Algebras}

In this chapter,† the letter K denotes a non-zero commutative ring. The unit element of K is denoted by 1. Unless otherwise mentioned, all cogebras, algebras and bigebras, all modules and all tensor products are over K.

From § 6, K will be assumed to be a field of characteristic 0.

\section{ENVELOPING BIGEBRA OF A LIE ALGEBRA}

Throughout this paragraph, \(\mathfrak{g}\) will denote a Lie algebra over \(\mathbf{K}\) , \(\mathrm{U}(\mathfrak{g})\) or simply \(\mathrm{U}\) its enveloping algebra (Chapter I, § 2, no. 1), \(\sigma\) the canonical mapping of \(\mathfrak{g}\) into \(\mathrm{U}(\mathfrak{g})\) (loc. cit.) and \((\mathrm{U}_n)_{n\geqslant 0}\) the canonical filtration of \(\mathrm{U}\) (loc. cit., no. 6).

\subsection{PRIMITIVE ELEMENTS OF A COGEBRA}

Throughout this no. we consider a cogebra E (Algebra, Chapter III, § 11, no. 1) with coproduct

\[
c \colon \mathrm{E} \rightarrow \mathrm{E} \otimes \mathrm{E}
\]

and counit \(\varepsilon\) (loc. cit., no. 2). Recall that \(\varepsilon\) is a linear form on the K-module E such that (with the canonical identification of \(\mathrm{E} \otimes \mathrm{K}\) and \(\mathrm{K} \otimes \mathrm{E}\) with E):

\[
\mathrm{Id} _ {\mathrm{E}} = (\varepsilon \otimes \mathrm{Id} _ {\mathrm{E}}) \circ c = (\mathrm{Id} _ {\mathrm{E}} \otimes \varepsilon) \circ c.
\]

Let \(E^{+}\) denote the kernel of \(\varepsilon\) and let u be an element of E such that

\[
c (u) = u \otimes u \quad \text { and } \quad \varepsilon (u) = 1.
\]

The K-module E is the direct sum of \(E^{+}\) and the submodule K.u which is free with basis u; let \(\pi_{u}: E \to E^{+}\) and \(\eta_{u}: E \to K.u\) denote the projectors associated with this decomposition. Then

\[
\pi_ {u} (x) = x - \varepsilon (x). u, \quad \eta_ {u} (x) = \varepsilon (x). u.\tag{1}
\]

DEFINITION 1. An element \(x\) of \(\mathbf{E}\) is called \(u\) -primitive if

\[
c (x) = x \otimes u + u \otimes x.\tag{2}
\]

The \(u\) -primitive elements of \(\mathbf{E}\) form a submodule of \(\mathbf{E}\) , denoted by \(\mathrm{P}_{u}(\mathrm{E})\) .

PROPOSITION 1. Every \(u\) -primitive element of \(\mathbf{E}\) belongs to \(\mathbf{E}^{+}\) .

\begin{enumerate}
\def\labelenumi{(\arabic{enumi})}
\setcounter{enumi}{1}
\tightlist
\item
  implies \(x = \varepsilon(x).u + \varepsilon(u).x = \varepsilon(x).u + x\) , whence \(\varepsilon(x) = 0\) .
\end{enumerate}

Remark. If \(x \in \mathrm{E}\) and \(c(x) = x' \otimes u + u \otimes x''\) , where \(x', x''\) are in \(\mathrm{E}^+\) , then \(x = \varepsilon(x')\) . \(u + \varepsilon(u)\) . \(x'' = x''\) ; similarly \(x = x'\) and \(x\) is \(u\) -primitive.

For all \(x \in \mathrm{E}^{+}\) , we write

\[
c _ {u} ^ {+} (x) = c (x) - x \otimes u - u \otimes x.\tag{3}
\]

PROPOSITION 2. We have

\[
\left(\pi_ {u} \otimes \pi_ {u}\right) \circ c = c _ {u} ^ {+} \circ \pi_ {u}.\tag{4}
\]

Let \(x\) be in E; then

\[
\begin{array}{r l} (\pi_ {u} \otimes \pi_ {u}) (c (x)) & = ((1 - \eta_ {u}) \otimes (1 - \eta_ {u})) (c (x)) \\ & = c (x) - (1 \otimes \eta_ {u}) (c (x)) - (\eta_ {u} \otimes 1) (c (x)) + (\eta_ {u} \otimes \eta_ {u}) (c (x)). \end{array}
\]

As \(\varepsilon\) is counit of E,

\[
(1 \otimes \eta_ {u}) (c (x)) = x \otimes u, \quad (\eta_ {u} \otimes 1) (c (x)) = u \otimes x
\]

whence

\[
\left(\eta_ {u} \otimes \eta_ {u}\right) (c (x)) = \left(\eta_ {u} \otimes 1\right) \left(\left(1 \otimes \eta_ {u}\right) (c (x))\right) = \varepsilon (x). u \otimes u;
\]

from this we conclude

\[
\left(\pi_ {u} \otimes \pi_ {u}\right) (c (x)) = c (x) - x \otimes u - u \otimes x + \varepsilon (x). u \otimes u.
\]

On the other hand,

\[
c _ {u} ^ {+} \left(\pi_ {u} (x)\right) = c (x) - x \otimes u - u \otimes x + \varepsilon (x). u \otimes u,
\]

whence formula (4).

As \(\mathrm{E}^{+}\) is a direct factor submodule of \(\mathrm{E}\) , \(\mathrm{E}^{+} \otimes \mathrm{E}^{+}\) can be identified with a direct factor submodule of \(\mathrm{E} \otimes \mathrm{E}\) . With this identification, \(\pi_{u} \otimes \pi_{u}\) is a projector of \(\mathrm{E} \otimes \mathrm{E}\) onto \(\mathrm{E}^{+} \otimes \mathrm{E}^{+}\) . By formula (4), \(c_{u}^{+}\) maps \(\mathrm{E}^{+}\) into \(\mathrm{E}^{+} \otimes \mathrm{E}^{+}\) and \(\pi_{u}\) is a morphism of the cogebra \((\mathrm{E}, c)\) into the cogebra \((\mathrm{E}^{+}, c_{u}^{+})\) .

PROPOSITION 3. If the cogebra \((\mathrm{E}, c)\) is coassociative (resp. cocommutative) (Algebra, Chapter III, § 11, no. 2), so is the cogebra \((\mathrm{E}^{+}, c_{\mathbf{u}}^{+})\) . This follows from the following lemma:

Lemma 1. Let \(\pi: \mathrm{E} \to \mathrm{E}'\) be a surjective cogebra morphism. If \(\mathrm{E}\) is coassociative (resp. cocommutative), so is \(\mathrm{E}'\) .

Let B be an associative K-algebra; the mapping \(f \mapsto f \circ \pi\) is an injective algebra homomorphism of \(\mathrm{Hom}_{\mathbb{K}}(\mathrm{E}', \mathrm{B})\) into \(\mathrm{Hom}_{\mathbb{K}}(\mathrm{E}, \mathrm{B})\) . It then suffices to apply Proposition 1 (resp. Proposition 2) of Algebra, Chapter III, § 11, no. 2.

\subsection{PRIMITIVE ELEMENTS OF A BIGEBRA}

Let E be a bigebra (Algebra, Chapter III, § 11, no. 4), c its coproduct, ε its counit and 1 its unit element. As ε(1) = 1 and c(1) = 1 ⊗ 1, the results of the previous no. can be applied with u = 1. The 1-primitive elements of E (no. 1, Definition 1) are simply called primitive (cf.~Algebra, Chapter III; § 11, no. 8), that is the elements x of E such that

\[
c (x) = x \otimes 1 + 1 \otimes x.\tag{5}
\]

We write simply \(\pi\) , \(\eta\) , P(E), \(c^{+}\) instead of \(\pi_{1}\) , \(\eta_{1}\) , P \(_{1}\) (E), \(c_{1}^{+}\) .

PROPOSITION 4. The set \(\mathrm{P(E)}\) of primitive elements of \(\mathbf{E}\) is a Lie subalgebra of \(\mathbf{E}\) .

If \(x, y\) are in \(\mathrm{P(E)}\) , then

\[
\begin{array}{r l} c (x y) & = c (x) c (y) = (x \otimes 1 + 1 \otimes x) (y \otimes 1 + 1 \otimes y) \\ & = x y \otimes 1 + 1 \otimes x y + x \otimes y + y \otimes x, \end{array}
\]

whence

\[
c ([ x, y ]) = [ x, y ] \otimes 1 + 1 \otimes [ x, y ].
\]

PROPOSITION 5. Let \(f: \mathrm{E} \to \mathrm{E}'\) be a bigebra morphism. If \(x\) is a primitive element of \(\mathrm{E}\) , then \(f(x)\) is a primitive element of \(\mathrm{E}'\) and the restriction of \(f\) to \(\mathrm{P(E)}\) is a Lie algebra homomorphism \(\mathrm{P}(f): \mathrm{P(E)} \to \mathrm{P(E')}\) .

Let \(c\) (resp. \(c'\) ) be the coproduct of E (resp. E'). Since \(f\) is a cogebra morphism \(c' \circ f = (f \otimes f) \circ c\) , whence

\[
\begin{array}{r l} c ^ {\prime} (f (x)) & = (f \otimes f) (c (x)) = (f \otimes f) (x \otimes 1 + 1 \otimes x) \\ & = f (x) \otimes 1 + 1 \otimes f (x), \end{array}
\]

for \(x\) primitive. Hence \(f\) maps \(\mathrm{P(E)}\) into \(\mathrm{P(E')}\) and \(f([x,y]) = [f(x),f(y)]\) since \(f\) is an algebra homomorphism.

Remarks. (1) Let \(p\) be a prime number such that \(p\) . \(1 = 0\) in \(\mathbf{K}\) . The binomial formula and the congruences \(\binom{p}{i} \equiv 0 \pmod{p}\) for \(1 \leqslant i \leqslant p - 1\) imply that

P(E) is stable under the mapping \(x \mapsto x^p\) .

\begin{enumerate}
\def\labelenumi{(\arabic{enumi})}
\setcounter{enumi}{1}
\tightlist
\item
  By definition, the diagram
\end{enumerate}

\[
0 \longrightarrow \mathrm{P(E)} \longrightarrow \mathrm {E^ {+}} \stackrel {c ^ {+}} {\longrightarrow} \mathrm {E^ {+}} \otimes \mathrm {E^ {+}}
\]

is an exact sequence. If \(\mathbf{K}'\) is a commutative ring and \(\rho: \mathbf{K} \to \mathbf{K}'\) a ring homomorphism, \(\rho^*(\mathrm{E}) = \mathrm{E} \otimes_{\mathbb{K}} \mathbf{K}'\) is a \(\mathbf{K}'\) -bigebra and the inclusion \(\mathrm{P}(\mathrm{E}) \to \mathrm{E}\) defines a homomorphism of Lie \(\mathbf{K}'\) -algebras

\[
\alpha \colon \mathrm{P} (\mathrm{E}) \otimes_ {\mathrm{K}} \mathrm{K} ^ {\prime} \rightarrow \mathrm{P} (\mathrm{E} \otimes_ {\mathrm{K}} \mathrm{K} ^ {\prime}).
\]

If \(\mathbf{K}'\) is flat over \(\mathbf{K}\) (Commutative Algebra, Chapter I, §2, no. 3, Definition 2), it follows from loc. cit. that the diagram

\[
0 \longrightarrow \mathrm{P} (\mathrm{E}) \otimes_ {\mathrm{K}} \mathrm{K} ^ {\prime} \longrightarrow \mathrm{E} ^ {+} \otimes_ {\mathrm{K}} \mathrm{K} ^ {\prime} \xrightarrow {c ^ {+} \otimes_ {\mathrm{K}} \mathrm{Id} _ {\mathrm{K} ^ {\prime}}} \left(\mathrm{E} ^ {+} \otimes_ {\mathrm{K}} \mathrm{K} ^ {\prime}\right) \otimes_ {\mathrm{K} ^ {\prime}} \left(\mathrm{E} ^ {+} \otimes_ {\mathrm{K}} \mathrm{K} ^ {\prime}\right)
\]

is an exact sequence, which implies that \(\alpha\) is an isomorphism.

\subsection{FILTERED BIGEBRAS}

DEFINITION 2. Let E be a bigebra with coproduct c.~A filtration compatible with the bigebra structure on E is an increasing sequence \((\mathrm{E}_{n})_{n\geq0}\) of submodules of E such that

\[
\begin{array}{c} \mathrm{E} _ {0} = \mathrm{K}. 1, \qquad \mathrm{E} = \bigcup_ {n \geqslant 0} \mathrm{E} _ {n} \\ \mathrm{E} _ {m}. \mathrm{E} _ {n} \subset \mathrm{E} _ {m + n} \quad \text {for} \quad m \geqslant 0, n \geqslant 0 \\ c (\mathrm{E} _ {n}) \subset \sum_ {i + j = n} \mathrm{Im} (\mathrm{E} _ {i} \otimes \mathrm{E} _ {j}) \quad \text {for} n \geqslant 0. \dagger \end{array}\tag{6}
\]

A bigebra with a filtration compatible with its bigebra structure is called a filtered bigebra.

Example. Let E be a graded bigebra (Algebra, Chapter III, § 11, no. 4, Definition 3) and \((\mathrm{E}^n)_{n\geqslant 0}\) its graduation. We write \(\mathrm{E}_n = \sum_{i = 0}^{n}\mathrm{E}^i\) . The sequence \((\mathrm{E}_n)\) is a filtration compatible with the bigebra structure on E.

PROPOSITION 6. Let E be a filtered bigebra and \((\mathrm{E}_n)_{n\geqslant 0}\) its filtration. For every integer \(n\geqslant 0\) , let \(\mathrm{E}_n^+ = \mathrm{E}_n\cap \mathrm{E}^+\) . Then \(\mathrm{E_0^+} = \{0\}\) and

\[
c ^ {+} (\mathrm{E} _ {n} ^ {+}) \subset \sum_ {i = 1} ^ {n - 1} \operatorname{Im} (\mathrm{E} _ {i} ^ {+} \otimes \mathrm{E} _ {n - i} ^ {+}) \quad \text { for } n \geqslant 0. \dagger\tag{7}
\]

As \(\mathrm{E}_0 = \mathrm{K}.1, \mathrm{E}_0^+ = 0\) . If \(x \in \mathrm{E}_n\) , \(\pi(x) = x - \varepsilon(x).1\) (formula (1)), whence \(\pi(x) \in \mathrm{E}_n^+\) and \(\pi(\mathrm{E}_n) \subset \mathrm{E}_n^+\) . It follows that \(\pi \otimes \pi\) maps \(\operatorname{Im}(\mathrm{E}_i \otimes \mathrm{E}_j)\) into \(\operatorname{Im}(\mathrm{E}_i^+ \otimes \mathrm{E}_j^+)\) for \(i \geqslant 0, j \geqslant 0\) . As \(c^+ = (\pi \otimes \pi) \circ c\) in \(\mathrm{E}^+\) (no. 1, Proposition 2), by (6)

\[
c ^ {+} \left(\mathrm{E} _ {n} ^ {+}\right) \subset \sum_ {i = 0} ^ {n} \operatorname{Im} \left(\mathrm{E} _ {i} ^ {+} \otimes \mathrm{E} _ {n - i} ^ {+}\right) = \sum_ {i = 1} ^ {n - 1} \operatorname{Im} \left(\mathrm{E} _ {i} ^ {+} \otimes \mathrm{E} _ {n - i} ^ {+}\right).
\]

COROLLARY. The elements of \(\mathbf{E}_1^+\) are primitive.

If \(x \in \mathrm{E}_{1}^{+}\) , then \(c^{+}(x) = 0\) by (7), whence (5).

\subsection{ENVELOPING BIGEBRA OF A LIE ALGEBRA}

Recall that g denotes a Lie algebra and U its enveloping algebra, with its canonical filtration \((\mathrm{U}_{n})_{n\geqslant0}\) .

PROPOSITION 7. There exists on the algebra U one and only one coproduct c which makes U into a bigebra such that the elements of \(\sigma(g)\) are primitive. The bigebra (U, c) is cocommutative; its counit is the linear form \(\varepsilon\) such that the constant term (Chapter I, §2, no. 1) of every element \(x\) of U is \(\varepsilon(x)\) .1. The canonical filtration \((\mathrm{U}_n)_{n\geqslant 0}\) of U is compatible with this bigebra structure.

\begin{enumerate}
\def\labelenumi{(\alph{enumi})}
\tightlist
\item
  Let \(x \in \mathfrak{g}\) ; we write \(c_0(x) = \sigma(x) \otimes 1 + 1 \otimes \sigma(x) \in \mathrm{U} \otimes \mathrm{U}\) . If \(x, y\) are in \(\mathfrak{g}\) , then
\end{enumerate}

\[
c _ {0} (x) c _ {0} (y) = (\sigma (x) \sigma (y)) \otimes 1 + 1 \otimes (\sigma (x) \sigma (y)) + \sigma (x) \otimes \sigma (y) + \sigma (y) \otimes \sigma (x),
\]

whence

\[
[ c _ {0} (x), c _ {0} (y) ] = c _ {0} ([ x, y ]).
\]

By the universal property of U (Chapter I, §2, no. 1, Proposition 1), there exists one and only one unital algebra homomorphism

\[
c \colon \mathrm{U} \to \mathrm{U} \otimes \mathrm{U}
\]

such that \(c(\sigma(x)) = \sigma(x) \otimes 1 + 1 \otimes \sigma(x)\) for \(x \in \mathfrak{g}\) . This proves the uniqueness assertion of Proposition 7.

\begin{enumerate}
\def\labelenumi{(\alph{enumi})}
\setcounter{enumi}{1}
\tightlist
\item
  We show that \(c\) is coassociative. The linear mappings \(c'\) and \(c''\) of \(U\) into \(U \otimes U \otimes U\) defined by
\end{enumerate}

\[
c ^ {\prime} = (c \otimes \mathrm{Id} _ {\mathrm{U}}) \circ c \quad \text { and } \quad c ^ {\prime \prime} = (\mathrm{Id} _ {\mathrm{U}} \otimes c) \circ c
\]

are unital algebra homomorphisms which coincide on \(\sigma (\mathfrak{g})\) since, for \(a\in \sigma (\mathfrak{g})\)

\[
c ^ {\prime} (a) = a \otimes 1 \otimes 1 + 1 \otimes a \otimes 1 + 1 \otimes 1 \otimes a = c ^ {\prime \prime} (a),
\]

whence the result.

\begin{enumerate}
\def\labelenumi{(\alph{enumi})}
\setcounter{enumi}{2}
\item
  We show that \(c\) is cocommutative. Let \(\tau\) be the automorphism of \(\mathbf{U} \otimes \mathbf{U}\) such that \(\tau(a \otimes b) = b \otimes a\) for \(a, b\) in \(\mathbf{U}\) . The mappings \(\tau \circ c\) and \(c\) of \(\mathbf{U}\) into \(\mathbf{U} \otimes \mathbf{U}\) are unital algebra homomorphisms which coincide on \(\sigma(\mathfrak{g})\) , whence the result.
\item
  We show that \(\varepsilon\) is a counit for \(c\) . The mappings \((\mathrm{Id}_{\mathbb{U}} \otimes \varepsilon) \circ c\) and \((\varepsilon \otimes \mathrm{Id}_{\mathbb{U}}) \circ c\) , of U into U are unital algebra homomorphisms which coincide with \(\mathrm{Id}_{\mathbb{U}}\) on \(\sigma(\mathfrak{g})\) .
\item
  We know that \(\mathrm{U}_0 = \mathrm{K}.1, \mathrm{U}_n \subset \mathrm{U}_{n+1}, \mathrm{U} = \bigcup_{n \geqslant 0} \mathrm{U}_n\) and \(\mathrm{U}_n. \mathrm{U}_m \subset \mathrm{U}_{n+m}\) (Chapter I, §2, no. 6). Let \(a_1, \ldots, a_n\) be in \(\sigma(\mathfrak{g})\) . Then
\end{enumerate}

\[
\begin{array}{l} c (a _ {1} \ldots a _ {n}) = \prod_ {i = 1} ^ {n} c (a _ {i}) = \prod_ {i = 1} ^ {n} (a _ {i} \otimes 1 + 1 \otimes a _ {i}) \\ = \sum_ {i = 0} ^ {n} \sum_ {\alpha \in I (i)} (a _ {\alpha (1)} \ldots a _ {\alpha (i)}) \otimes (a _ {\alpha (i + 1)} \ldots a _ {\alpha (n)}), \end{array}\tag{8}
\]

where \(I(i)\) denotes the set of permutations of \([1, n]\) which are increasing in each of the intervals \([1, i]\) and \([i + 1, n]\) . As \(U_{n}\) is the K-module generated by the products of at most n elements of \(\sigma(g)\) , formula (8) implies that the filtration \((\mathrm{U}_{n})\) is compatible with the bigebra structure of \((\mathrm{U}, c)\) .

DEFINITION 3. The bigebra \((\mathrm{U}, c)\) is called the enveloping bigebra of the Lie algebra \(\mathfrak{g}\) .

PROPOSITION 8. Let E be a bigebra with coproduct denoted by \(c_{\mathbb{E}}\) and let h be a Lie algebra homomorphism of g into P(E) (no. 2, Proposition 4). The unital algebra homomorphism f: U → E such that f(σ(x)) = h(x) for all x ∈ g is a bigebra morphism.

\[
(f \otimes f) \circ c = c _ {\mathrm{E}} \circ f.
\]

\[
a \in \sigma (\mathfrak {g})
\]

\[
(f \otimes f) (c (a)) = f (a) \otimes 1 + 1 \otimes f (a) = c _ {\mathbb {E}} (f (a))
\]

since \(f(a) \in \mathrm{P}(\mathrm{E})\) . Similarly if \(\varepsilon_{\mathbb{E}}\) is the counit of \(\mathrm{E}\) , \(\varepsilon_{\mathbb{E}} \circ f\) is a unital algebra homomorphism \(\mathrm{U} \to \mathrm{K}\) which is zero on \(\sigma(\mathfrak{g})\) (no. 1, Proposition 1) and therefore coincides with \(\varepsilon\) .

It follows from Propositions 5 and 8 that the mapping \(f \mapsto f \circ \sigma\) defines a one-to-one correspondence between bigebra homomorphisms \(U(\mathfrak{g}) \to E\) and Lie algebra homomorphisms \(\mathfrak{g} \to P(E)\) .

COROLLARY. Let \(\mathfrak{g}_i\) ( \(i = 1,2\) ) be a Lie algebra, \(\mathrm{U}(\mathfrak{g}_i)\) its enveloping bigebra and \(\sigma_i\colon \mathfrak{g}_i\to \mathrm{U}(\mathfrak{g}_i)\) the canonical mapping. For every Lie algebra homomorphism \(h\colon \mathfrak{g}_1\to \mathfrak{g}_2\) , the unital algebra homomorphism \(\mathrm{U}(h)\colon \mathrm{U}(\mathfrak{g}_1)\to \mathrm{U}(\mathfrak{g}_2)\) such that \(\mathrm{U}(h)\circ \sigma_1 = \sigma_2\circ h\) (Chapter I, § 2, no. 1) is a bigebra morphism.

\subsection{STRUCTURE OF THE COGEBRA U(g) IN CHARACTERISTIC 0}

In this no., K will be assumed to be a field of characteristic 0.

Let \(S(\mathfrak{g})\) be the symmetric algebra of the vector space \(\mathfrak{g}\) , \(c_{\mathrm{s}}\) its coproduct (Algebra, Chapter III, § 11, no. 1, Example 6) and \(\eta\) the canonical isomorphism of the vector space \(\mathsf{S}(\mathfrak{g})\) onto the vector space U (Chapter I, § 2, no. 7). Recall that if \(x_{1},\ldots ,x_{n}\) are in \(\mathfrak{g}\) , then

\[
\eta (x _ {1} \dots x _ {n}) = \frac {1}{n !} \sum_ {\tau \in \mathfrak {S} _ {n}} \sigma (x _ {\tau (1)}) \dots \sigma (x _ {\tau (n)}).\tag{9}
\]

In particular, for \(x \in g\) and \(n \geqslant 0\) ,

\[
\eta (x ^ {n}) = \sigma (x) ^ {n}.\tag{10}
\]

Note that by Algebra, Chapter III, § 6, no. 1, Remark 3, \(\eta\) is the unique linear mapping of \(\mathbb{S}(\mathfrak{g})\) into U satisfying condition (10).

PROPOSITION 9. For every integer \(n \geqslant 0\) , let \(\mathbf{U}^n\) be the vector subspace of \(\mathbf{U}\) generated by the \(\sigma(x)^n\) for \(x \in \mathfrak{g}\) .

\begin{enumerate}
\def\labelenumi{(\alph{enumi})}
\tightlist
\item
  The sequence \((\mathbf{U}^n)_{n\geqslant 0}\) is a graduation of the vector space \(\mathbf{U}\) compatible with its cogebra structure.
\end{enumerate}

Let U be given the graduation (U \(^{n}\) ).

\begin{enumerate}
\def\labelenumi{(\alph{enumi})}
\setcounter{enumi}{1}
\tightlist
\item
  The canonical mapping \(\eta: \mathbf{S}(\mathfrak{g}) \to \mathbf{U}\) is an isomorphism of graded cogebras. Let \(x \in \mathfrak{g}\) and \(n \in \mathbf{N}\) . Then
\end{enumerate}

\[
c _ {\mathrm{S}} (x ^ {n}) = c _ {\mathrm{S}} (x) ^ {n} = (x \otimes 1 + 1 \otimes x) ^ {n} = \sum_ {i = 0} ^ {n} {\binom {n} {i}} x ^ {i} \otimes x ^ {n - i}\tag{11}
\]

since \(c_{\mathbb{S}}\) is an algebra homomorphism. Similarly, by (10),

\[
\begin{array}{l} c (\eta (x ^ {n})) = c (\sigma (x) ^ {n}) = c (\sigma (x)) ^ {n} = (\sigma (x) \otimes 1 + 1 \otimes \sigma (x)) ^ {n} \\ = \sum_ {i = 0} ^ {n} \binom {n} {i} \sigma (x) ^ {i} \otimes \sigma (x) ^ {n - i} = \sum_ {i = 0} ^ {n} \binom {n} {i} \eta (x ^ {i}) \otimes \eta (x ^ {n - i}), \end{array}\tag{12}
\]

whence

\[
(\eta \otimes \eta) (c _ {\mathrm{S}} (x ^ {n})) = c (\eta (x ^ {n})).
\]

As the \(x^n\) , for \(x \in \mathfrak{g}\) and \(n \in \mathbf{N}\) , generate the vector space \(\mathsf{S}(\mathfrak{g})\) , \((\eta \otimes \eta) \circ c_{\mathbf{S}} = c \circ \eta\) and \(\eta\) is a cogebra isomorphism.

On the other hand, formula (10) shows that \(\eta(\mathbb{S}^n(\mathfrak{g})) = \mathrm{U}^n\) , which completes the proof of (a) and (b) taking account of the fact that the graduation of \(\mathbb{S}(\mathfrak{g})\) is compatible with its cogebra structure.

The graduation \((\mathrm{U}^{n})_{n\geqslant0}\) of U is called the canonical graduation.

COROLLARY. The canonical mapping \(\sigma\) defines an isomorphism of \(g\) onto the Lie algebra \(\mathrm{P}(\mathbf{U})\) of primitive elements of \(\mathbf{U}\) .

As \(c^{+}\) is a graded homomorphism of degree 0,

\[
\mathrm{P} (\mathrm{U}) = \sum_ {n \geqslant 1} (\mathrm{P} (\mathrm{U}) \cap \mathrm{U} ^ {n}).
\]

It suffices to prove that if \(n > 1\) and \(a \in \mathrm{U}^n\) is primitive, then \(a = 0\) . Now \(a\) can be written as \(\sum_{i} \lambda_i a_i^n\) , where \(\lambda_i \in \mathbf{K}\) , \(a_i \in \sigma(g)\) . By (12), the term of bi-degree \((1, n - 1)\) in \(c^+(a)\) is \(n \sum_{i} \lambda_i a_i \otimes a_i^{n-1}\) . Hence \(\sum_{i} \lambda_i a_i \otimes a_i^{n-1} = 0\) . If \(\mu: \mathrm{U} \otimes \mathrm{U} \to \mathrm{U}\) is the linear mapping defined by multiplication on \(\mathrm{U}\) , then

\[
a = \sum_ {i} \lambda_ {i} a _ {i} ^ {n} = \mu \left(\sum_ {i} \lambda_ {i} a _ {i} \otimes a _ {i} ^ {n - 1}\right) = 0.
\]

Remarks. (1) \(\mathrm{U}_n = \sum_{i=0}^{n} \mathrm{U}^i\) (Chapter I, § 2, no. 7, Corollary 4 to Theorem 1).

\begin{enumerate}
\def\labelenumi{(\arabic{enumi})}
\setcounter{enumi}{1}
\tightlist
\item
  The mapping \(\eta\) is the unique morphism of graded cogebras of \(\mathsf{S}(\mathfrak{g})\) into U such that \(\eta(1) = 1\) and \(\eta(x) = \sigma(x)\) for \(x \in \mathfrak{g}\) . For if \(\eta'\) is a morphism satisfying these conditions, we prove by induction on \(n\) that \(\eta'(x^n) = \eta(x^n)\) for \(x \in \mathfrak{g}\) and \(n > 1\) . As \(c_{\mathbb{S}}^{+}(x^{n}) = \sum_{i=1}^{n-1} \binom{n}{i} x^{i} \otimes x^{n-i}\) by (3) and (11),
\end{enumerate}

\[
(\eta \otimes \eta) (c _ {\mathbf {S}} ^ {+} (x ^ {n})) = (\eta^ {\prime} \otimes \eta^ {\prime}) (c _ {\mathbf {S}} ^ {+} (x ^ {n}))
\]

by the induction hypothesis. It follows that \(c^{+}(\eta(x^{n})) = c^{+}(\eta'(x^{n}))\) ; it follows that \(\eta(x^{n}) - \eta'(x^{n})\) is a primitive element of degree \(n\) and hence is zero (Corollary to Proposition 9).

\begin{enumerate}
\def\labelenumi{(\arabic{enumi})}
\setcounter{enumi}{2}
\tightlist
\item
  Let \(\psi\) be the canonical isomorphism of the bigebra \(\mathsf{TS}(\mathfrak{g})\) onto the bigebra \(\mathsf{S}(\mathfrak{g})\) (Algebra, Chapter IV, § 5, Corollary 1 to Proposition 12). The mapping
\end{enumerate}

\[
\eta \circ \psi : T S (\mathfrak {g}) \rightarrow U
\]

is called canonical. It is the unique morphism \(\eta'\) of graded cogebras of TS(g) into U such that \(\eta'(1) = 1\) and \(\eta'(x) = \sigma(x)\) for all \(x \in \mathfrak{g}\) .

\begin{enumerate}
\def\labelenumi{(\arabic{enumi})}
\setcounter{enumi}{3}
\tightlist
\item
  Let V be a vector space. The primitive elements of the bigebra \(S(V)\) are the elements of degree 1. This follows from the Corollary to Proposition 9 applied to the commutative Lie algebra V.
\end{enumerate}

Let \((e_i)_{i\in I}\) be a basis of the vector K-space \(g\) , where the indexing set I is totally ordered. For all \(\alpha \in \mathbf{N}^{(I)}\) , we write

\[
e _ {\alpha} = \prod_ {i \in I} \frac {\sigma (e _ {i}) ^ {\alpha (i)}}{\alpha (i) !}.\tag{13}
\]

The \(e_{\alpha}\) , for \(|\alpha| \leqslant n\) , form a basis of the vector K-space \(\mathbf{U}_n\) (Chapter I, §2, no. 7, Corollary 3 to Theorem 1). Then

\[
e _ {0} = 1, \quad e _ {\varepsilon_ {i}} = \sigma (e _ {i}) \quad \text { for } i \in I.
\]

As the graded algebra associated with the filtered algebra U is commutative (loc. cit., Theorem 1), for \(\alpha, \beta\) in \(\mathbf{N}^{(I)}\) ,

\[
e _ {\alpha}. e _ {\beta} \equiv ((\alpha , \beta)). e _ {\alpha + \beta} \mod U _ {| \alpha | + | \beta | - 1}.\tag{14}
\]

where

\[
((\alpha , \beta)) = \prod_ {i \in I} \frac {(\alpha (i) + \beta (i)) !}{\alpha (i) ! \beta (i) !}.
\]

On the other hand, we have immediately

\[
\varepsilon (e _ {0}) = 1, \quad \varepsilon (e _ {\alpha}) = 0 \quad \text { for } | \alpha | \geqslant 1.\tag{15}
\]

Finally, formula (12) implies that, for \(\alpha \in \mathbf{N}^{(I)}\) ,

\[
c (e _ {\alpha}) = \sum_ {\beta + \gamma = \alpha} e _ {\beta} \otimes e _ {\gamma}.\tag{16}
\]

This formula allows us to determine the algebra \(\mathrm{U}' = \mathrm{Hom}(\mathrm{U},\mathrm{K})\) dual to the cogebra U (Algebra, Chapter III, § 11, no. 2). For let \(\mathrm{K}[[\mathrm{X}_i]_{i\in \mathbf{I}}]\) be the algebra of formal power series in indeterminates \((\mathrm{X}_i)_{i\in \mathbf{I}}\) (cf.~Algebra, Chapter III, § 2, no. 11); if \(\lambda \in \mathrm{U}'\) , let \(f_{\lambda}\) denote the formal power series

\[
f _ {\lambda} = \sum_ {\alpha} \langle \lambda , e _ {\alpha} \rangle \mathrm{X} ^ {\alpha}, \text { where } \mathrm{X} ^ {\alpha} = \prod_ {i \in I} \mathrm{X} _ {i} ^ {\alpha (i)}
\]

and the summation index \(\alpha\) runs through \(\mathbf{N}^{(I)}\) .

PROPOSITION 10. The mapping \(\lambda \mapsto f_{\lambda}\) is an isomorphism of the algebra \(\mathbf{U}'\) onto the algebra of formal power series \(\mathbf{K}[[\mathbf{X}_i]]_{i \in I}\) .

Because \((e_{\alpha})\) is a basis of U, the mapping \(\lambda \mapsto f_{\lambda}\) is K-linear and bijective. On the other hand, for \(\lambda, \mu\) in U',

\[
\begin{array}{r l} f _ {\lambda \mu} & = \sum_ {\alpha} \langle \lambda \mu , e _ {\alpha} \rangle \mathbf {X} ^ {\alpha} = \sum_ {\alpha} \langle \lambda \otimes \mu , c (e _ {\alpha}) \rangle \mathbf {X} ^ {\alpha} \\ & = \sum_ {\alpha} \langle \lambda \otimes \mu , \sum_ {\beta + \gamma = \alpha} e _ {\beta} \otimes e _ {\gamma} \rangle \mathbf {X} ^ {\alpha} \\ & = \sum_ {\beta , \gamma} \langle \lambda , e _ {\beta} \rangle \langle \mu , e _ {\gamma} \rangle \mathbf {X} ^ {\beta + \gamma} = f _ {\lambda} f _ {\mu}, \end{array}\tag{by (16)}
\]

which shows that \(\lambda \mapsto f_{\lambda}\) is an algebra isomorphism and completes the proof.

\subsection{STRUCTURE OF FILTERED BIGEBRAS IN CHARACTERISTIC 0}

In this no. we continue to assume that K is a field of characteristic 0.

If E is a bigebra, the canonical injection \(P(E) \rightarrow E\) can be extended to a bigebra morphism \(f_{E}: U(P(E)) \rightarrow E\) (no. 4, Proposition 3).

THEOREM 1. Let E be a cocommutative bigebra.

\begin{enumerate}
\def\labelenumi{(\alph{enumi})}
\item
  The bigebra morphism \(f_{\mathbf{E}} \colon \mathrm{U}(\mathrm{P}(\mathbf{E})) \to \mathbf{E}\) is injective.
\item
  If there exists on E a filtration compatible with its bigebra structure (no. 3, Definition 2), the morphism \(f_{\mathrm{E}}\) is an isomorphism.
\end{enumerate}

(In case (b), the bigebra E is therefore identified with the enveloping bigebra of the Lie algebra of its primitive elements.)

Let \(c_{\mathbb{E}}\) (resp. \(\varepsilon_{\mathbb{E}}\) ) be the coproduct (resp. counit) of E. We write \(\mathfrak{g} = \mathrm{P}(\mathrm{E})\) ; let \((e_i)_{i \in \mathrm{I}}\) be a basis of the vector K-space \(\mathfrak{g}\) , where the indexing set I is totally ordered, and let \((e_\alpha)_{\alpha \in \mathbf{N}^{(I)}}\) be the basis introduced in the preceding no. We write \(\mathbf{X}_\alpha = f_{\mathbb{E}}(e_\alpha)\) for \(\alpha \in \mathbf{N}^{(I)}\) . By (15) and (16), we have:

\begin{enumerate}
\def\labelenumi{(\arabic{enumi})}
\setcounter{enumi}{16}
\tightlist
\item
\end{enumerate}

\[
\varepsilon_ {E} (X _ {0}) = 1, \quad \varepsilon_ {E} (X _ {\alpha}) = 0 \quad \text { for } | \alpha | \geqslant 1,\tag{18}
\]

\[
c _ {E} (X _ {\alpha}) = \sum_ {\beta + \gamma = \alpha} X _ {\beta} \otimes X _ {\gamma} \quad \text { for } \alpha \in \mathbf {N} ^ {(I)},
\]

since \(f_{E}\) is a cogebra morphism.

We show that \(f_{E}\) is injective. This results from the following lemma:

Lemma 2. Let V be a vector space, E a cogebra and \(f\colon\mathsf{S}(\mathrm{V})\to\mathrm{E}\) a cogebra morphism. If the restriction of \(f\) to \(\mathsf{S}^{\circ}(\mathrm{V})+\mathsf{S}^{1}(\mathrm{V})\) is injective, then \(f\) is injective.

Let \(n \geqslant 0\) ; we write \(\mathbb{S}_n = \sum_{i \geqslant n} \mathbb{S}^i(\mathbf{V})\) and \(c_s\) for the coproduct of \(\mathbb{S}(\mathbf{V})\) and show by induction on \(n\) that \(f \mid \mathbb{S}_n\) is injective. Since the assertion is trivial for \(n = 0\) and \(n = 1\) , we assume that \(n \geqslant 2\) and let \(u \in \mathbb{S}_n\) be such that \(f(u) = 0\) . Then

\[
\begin{array}{r l} {0 = c _ {\mathbf {E}} (f (u))} & {= (f \otimes f) (c _ {\mathbf {s}} (u))} \\ & {= f (u) \otimes 1 + 1 \otimes f (u) + (f \otimes f) (c _ {\mathbf {s}} ^ {+} (u))} \\ & {= (f \otimes f) (c _ {\mathbf {s}} ^ {+} (u)).} \end{array}
\]

As \(c_{\mathbb{S}}^{+}(u) \in \mathbb{S}_{n-1} \otimes \mathbb{S}_{n-1}\) , by (11) the induction hypothesis shows that \(u\) is a primitive element of \(\mathbb{S}(\mathrm{V})\) , hence is of degree 1 (no. 5, Remark 4) and hence is zero, since \(f \mid \mathbb{S}^1(\mathrm{V})\) is injective.

It follows in particular that the family \((X_{\alpha})\) is free.

We show that \(f_{\mathbb{E}}\) is surjective if \(\mathrm{E}\) has a filtration compatible with its bigebra structure. Let \((\mathrm{E}_n)_{n\geqslant 0}\) be such a filtration and write \(\mathrm{E}_n^+ = \mathrm{E}_n\cap \mathrm{Ker}(\epsilon_{\mathbb{E}})\) . We show by induction on \(n\) that \(\mathrm{E}_n^+\) is contained in the image of \(f_{\mathbb{E}}\) . As \(\mathrm{E} = \mathrm{K}.1 + \bigcup_{n > 0}\mathrm{E}_n^+\) , this will imply the surjectivity of \(f_{\mathbb{E}}\) . The assertion is trivial for \(n = 0\) and follows from the Corollary to Proposition 6 of no. 3 for \(n = 1\) ; suppose henceforth that \(n\geqslant 2\) and let \(x\in \mathrm{E}_n^+\) . By Proposition 6 of no. 3,

\[
c _ {\mathrm{E}} ^ {+} (x) \in \sum_ {i = 1} ^ {n - 1} \mathrm{E} _ {i} ^ {+} \otimes \mathrm{E} _ {n - i} ^ {+}
\]

and by the induction hypothesis there exist scalars \(\lambda_{\alpha, \beta}\) , where \(\alpha, \beta\) are in \(\mathbf{N}^{(I)}\) , which are zero except for a finite number, such that

\[
c _ {\mathbb {E}} ^ {+} (x) = \sum_ {\alpha , \beta \neq 0} \lambda_ {\alpha , \beta} \mathrm{X} _ {\alpha} \otimes \mathrm{X} _ {\beta}.\tag{19}
\]

Hence by formula (18)

\[
(c _ {\mathrm{E}} ^ {+} \otimes \operatorname{Id} _ {\mathrm{E}}) (c _ {\mathrm{E}} ^ {+} (x)) = \sum_ {\alpha , \beta , \gamma \neq 0} \lambda_ {\alpha + \beta , \gamma} \mathbf {X} _ {\alpha} \otimes \mathbf {X} _ {\beta} \otimes \mathbf {X} _ {\gamma}
\]

\[
(\mathrm{Id} _ {\mathtt {E}} \otimes c _ {\mathtt {E}} ^ {+}) (c _ {\mathtt {E}} ^ {+} (x)) = \sum_ {\alpha , \beta , \gamma \neq 0} \lambda_ {\alpha , \beta + \gamma} \mathbf {X} _ {\alpha} \otimes \mathbf {X} _ {\beta} \otimes \mathbf {X} _ {\gamma}.
\]

By Proposition 3 of no. 1 and the linear independence of the \(X_{\alpha}\) it follows that

\[
\lambda_ {\alpha + \beta , \gamma} = \lambda_ {\alpha , \beta + \gamma} \quad \text { for } \alpha , \beta , \gamma \text { in } \mathbf {N} ^ {(I)} - \{0 \}.\tag{20}
\]

On the other hand, the coproduct \(c_{E}\) is cocommutative; the same argument as above implies

\[
\lambda_ {\alpha , \beta} = \lambda_ {\beta , \alpha} \quad \text { for } \alpha , \beta \text { in } \mathbf {N} ^ {(I)} - \{0 \}.\tag{21}
\]

Suppose that there exists a family of scalars \((\mu_{\alpha})\) with \(|\alpha| \geqslant 2\) , such that

\[
\mu_ {\alpha + \beta} = \lambda_ {\alpha , \beta} \quad \text { for } \alpha , \beta \text { in } \mathbf {N} ^ {(I)} - \{0 \}.\tag{22}
\]

Then

\[
c _ {\mathrm{E}} ^ {+} (x) = \sum_ {\alpha , \beta \neq 0} \mu_ {\alpha + \beta} \mathrm{X} _ {\alpha} \otimes \mathrm{X} _ {\beta} = \sum_ {| \gamma | \geqslant 2} \mu_ {\gamma} c _ {\mathrm{E}} ^ {+} (\mathrm{X} _ {\gamma})
\]

by formula (18), hence \(y = x - \sum_{|\gamma| \geq 2} \mu_\gamma X_\gamma\) is primitive and hence belongs to \(P(E) \subset \operatorname{Im}(f_E)\) . Hence

\[
x = y + \sum_ {| \gamma | \geqslant 2} \mu_ {\gamma} f _ {\mathrm{E}} (e _ {\gamma}) \in \operatorname{Im} (f _ {\mathrm{E}}).
\]

The proof will therefore be complete when we have proved the following lemma:

Lemma 3. If a family of scalars \((\lambda_{\alpha, \beta})\) of finite support (with \(\alpha, \beta\) in \(\mathbf{N}^{(I)} - \{0\}\) ) satisfies relations (20) and (21), there exists a family \((\mu_{\alpha})_{|\alpha| \geqslant 2}\) of finite support such that \(\mu_{\alpha + \beta} = \lambda_{\alpha, \beta}\) for \(\alpha, \beta\) non-zero.

It suffices to prove that

\[
\alpha + \beta = \gamma + \delta\tag{23}
\]

implies \(\lambda_{\alpha, \beta} = \lambda_{\gamma, \delta}\) for \(\alpha, \beta, \gamma, \delta\) non-zero. By Riesz's Decomposition Lemma (Algebra, Chapter VI, § 1, no. 10, Theorem 1) there exist \(\pi, \rho, \sigma, \tau\) in \(\mathbf{N}^{(I)}\) such that

\[
\alpha = \pi + \sigma , \quad \beta = \rho + \tau , \quad \gamma = \pi + \rho , \quad \delta = \sigma + \tau .
\]

Suppose \(\pi \neq 0\) ; as \(\sigma + \beta = \rho + \delta\) , relation (20) implies

\[
\lambda_ {\alpha , \beta} = \lambda_ {\pi + \sigma , \beta} = \lambda_ {\pi , \sigma + \delta} = \lambda_ {\pi , \rho + \delta} = \lambda_ {\pi + \rho , \delta} = \lambda_ {\gamma , \delta}.
\]

If on the other hand \(\pi = 0\) , then \(\beta = \gamma + \tau\) and \(\delta = \alpha + \tau\) , whence

\[
\lambda_ {\alpha , \beta} = \lambda_ {\alpha , \gamma + \tau} = \lambda_ {\alpha + \tau , \gamma} = \lambda_ {\delta , \gamma}
\]

by (20), but also \(\lambda_{\delta, \gamma} = \lambda_{\gamma, \delta}\) by (21), whence \(\lambda_{\alpha, \beta} = \lambda_{\gamma, \delta}\) .

\section{FREE LIE ALGEBRAS}

\subsection{REVISION OF FREE ALGEBRAS}

Let X be a set. Recall the construction of the free magma M(X) constructed on X (Algebra, Chapter I, § 7, no. 1). By induction on the integer \(n \geqslant 1\) , we define the sets \(\mathbf{X}_n\) by writing \(\mathbf{X}_1 = \mathbf{X}\) and taking \(\mathbf{X}_n\) to be the sum set of the sets \(\mathbf{X}_p \times \mathbf{X}_{n-p}\) where \(p = 1, 2, \ldots, n-1\) ; if X is finite, so is each \(\mathbf{X}_n\) . The sum set of the family \((\mathbf{X}_n)_{n \geqslant 1}\) is denoted by M(X); each of the sets \(\mathbf{X}_n\) (and in particular X) is identified with a subset of M(X). Let w and \(w'\) be in M(X); let p and q denote the integers such that \(w \in \mathbf{X}_p\) and \(w' \in \mathbf{X}_q\) and let \(n = p + q\) ; the image of the ordered pair \(\langle w, w' \rangle\) under the canonical injection of \(\mathbf{X}_p \times \mathbf{X}_{n-p}\) into \(\mathbf{X}_n\) is denoted by \(w.w'\) and called the product of w and \(w'\) . Every mapping of X into a magma M can be extended in a unique way to a magma homomorphism of M(X) into M.

Let \(w\) be in \(\mathrm{M}(\mathbf{X})\) ; the unique integer \(n\) such that \(w \in \mathbf{X}_n\) is called the length of \(w\) and denoted by \(l(w)\) . Then \(l(w.w') = l(w) + l(w')\) for \(w, w'\) in \(\mathrm{M}(\mathbf{X})\) . The set \(\mathbf{X}\) is the subset of \(\mathrm{M}(\mathbf{X})\) consisting of the elements of length 1. Every element \(w\) of length \(\geqslant 2\) can be written uniquely in the form \(w = w'.w''\) .

The algebra of the magma M(X) with coefficients in the ring K is denoted by \(\operatorname{Lib}(X)\) , or \(\operatorname{Lib}_{K}(X)\) when it is necessary to indicate the ring K. The set M(X) is a basis of the K-module \(\operatorname{Lib}(X)\) and X will therefore be identified with a subset of \(\operatorname{Lib}(X)\) . If A is an algebra, every mapping of X into A can be extended uniquely to a homomorphism of \(\operatorname{Lib}(X)\) into A (Algebra, Chapter III, § 2, no. 7, Proposition 7).

\subsection{CONSTRUCTION OF THE FREE LIE ALGEBRA}

DEFINITION 1. The free Lie algebra over the set \(\mathbf{X}\) is the quotient algebra

\[
\mathrm{L} (\mathrm{X}) = \operatorname{Lib} (\mathrm{X}) / \mathrm{a},
\]

where a is the two-sided ideal of \(\operatorname{Lib}(\mathbf{X})\) generated by the elements of one of the forms

\begin{enumerate}
\def\labelenumi{(\arabic{enumi})}
\tightlist
\item
\end{enumerate}

\[
\mathrm{Q} (a) = a. a \text {   for   } a \text {   in   } \operatorname{Lib} (\mathrm{X}),\tag{2}
\]

\[
J (a, b, c) = a. (b. c) + b. (c. a) + c. (a. b)
\]

for \(a, b, c\) in \(\operatorname{Lib}(\mathbf{X})\) .

Clearly \(\mathrm{L}(\mathrm{X})\) is a Lie K-algebra; the composition of two elements \(u, v\) of \(\mathrm{L}(\mathrm{X})\) will be denoted by \([u, v]\) . When it is necessary to indicate the ring \(\mathbf{K}\) , we write \(\mathrm{L}_{\mathbb{K}}(\mathbf{X})\) for \(\mathrm{L}(\mathbf{X})\) .

The following proposition justifies the name free Lie algebra given to L(X).

PROPOSITION 1. Let \(\psi\) be the canonical mapping of \(\operatorname{Lib}(\mathbf{X})\) onto \(\operatorname{L}(\mathbf{X})\) and \(\phi\) the restriction of \(\psi\) to \(\mathbf{X}\) . For every mapping \(f\) of \(\mathbf{X}\) into a Lie algebra \(\mathfrak{g}\) , there exists one and only one homomorphism \(\mathrm{F}:\mathrm{L}(\mathrm{X})\to \mathfrak{g}\) such that \(f = \mathrm{F}\circ \phi\) .

\begin{enumerate}
\def\labelenumi{(\alph{enumi})}
\item
  Existence of F: let h be the homomorphism of \(\operatorname{Lib}(\mathbf{X})\) into g extending \(f\) (no. 1) . For all a in \(\operatorname{Lib}(\mathbf{X})\) , \(h(\mathbf{Q}(a)) = h(a.a) = [h(a), h(a)] = 0\) ; similarly, the Jacobi identity satisfied by g implies that \(h(\mathbf{J}(a, b, c)) = 0\) for a, b, c in \(\operatorname{Lib}(\mathbf{X})\) . It follows that \(h(a) = 0\) , whence there is a homomorphism F of \(\operatorname{L}(\mathbf{X})\) into g such that \(h = F \circ \phi\) . By restricting to X, we obtain \(f = F \circ \phi\) .
\item
  Uniqueness of F: let \(\mathrm{F}^{\prime}:\mathrm{L}(\mathrm{X})\to\mathfrak{g}\) be a homomorphism such that \(f=\mathrm{F}^{\prime}\circ\phi\) . The homomorphisms \(F\circ\psi\) and \(F^{\prime}\circ\psi\) of \(\operatorname{Lib}(X)\) into g coincide on X and hence are equal; as \(\psi\) is surjective, \(F=F^{\prime}\) .
\end{enumerate}

COROLLARY 1. The family \((\phi (x))_{x\in \mathbf{X}}\) is free over \(\mathbf{K}\) in \(\mathrm{L}(\mathbf{X})\) .

Let \(x_{1}, x_{2}, \ldots, x_{n}\) be distinct elements of X and \(\lambda_{1}, \ldots, \lambda_{n}\) be elements of K such that

\[
\lambda_ {1}. \phi (x _ {1}) + \dots + \lambda_ {n}. \phi (x _ {n}) = 0.\tag{3}
\]

Let \(\mathfrak{g}\) be the commutative Lie algebra with \(\mathbf{K}\) as underlying module. For \(i = 1,2,\ldots ,n\) , there exists a homomorphism \(\mathrm{F}_i\) of \(\mathbf{L}(\mathbf{X})\) into \(\mathfrak{g}\) such that \(\mathrm{F_i}(\phi (x_i)) = 1\) and \(\mathrm{F_i}(\phi (x)) = 0\) for \(x\neq x_{i}\) (Proposition 1); applying \(\mathrm{F}_i\) to relation (3), we obtain \(\lambda_{i} = 0\) .

COROLLARY 2. Let \(\mathfrak{a}\) be a Lie algebra. Every extension of \(\mathrm{L}(\mathrm{X})\) by \(\mathfrak{a}\) is inessential.

Let \(\mathfrak{a} \xrightarrow{\lambda} \mathfrak{g} \xrightarrow{\mu} \mathrm{L}(\mathrm{X})\) be such an extension (Chapter I, § 1, no. 7). As \(\mu\) is surjective, there exists a mapping \(f\) of \(\mathrm{X}\) into \(\mathfrak{g}\) such that \(\phi = \mu \circ f\) . Let \(\mathrm{F}\) be the homomorphism of \(\mathrm{L}(\mathrm{X})\) into \(\mathfrak{g}\) such that \(f = \mathrm{F} \circ \phi\) (Proposition 1). Then \((\mu \circ \mathrm{F}) \circ \phi = \mu \circ f = \phi\) and Proposition 1 shows that \(\mu \circ \mathrm{F}\) is the identity automorphism of \(\mathrm{L}(\mathrm{X})\) . The given extension is therefore inessential (Chapter I, § 1, no. 7, Proposition 6 and Definition 6).

As the ring K is non-zero, Corollary 1 to Proposition 1 shows that \(\phi\) is injective. Hence the set X can be identified by means of \(\phi\) with its image in L(X); with this convention, X generates L(X) and every mapping of X into a Lie algebra \(\mathfrak{g}\) can be extended to a Lie algebra homomorphism of L(X) into \(\mathfrak{g}\) .

Remark. When X is empty, M(X) is empty and hence L(X) = \{0\}. If X consists of a single element x, the submodule K.x of L(X) is a Lie subalgebra of

\(L(X)\) ; as X generates \(L(X)\) , Corollary 1 to Proposition 1 shows that \(L(X)\) is a free module with basis \(\{x\}\) .

\subsection{PRESENTATIONS OF A LIE ALGEBRA}

Let \(\mathfrak{g}\) be a Lie algebra and \(\pmb{a} = (a_i)_{i\in I}\) a family of elements of \(\mathfrak{g}\) . Let \(f_{a}\) be the homomorphism of \(\mathrm{L(I)}\) into \(\mathfrak{g}\) mapping each \(i\in I\) to \(a_{i}\) . The image of \(f_{a}\) is the subalgebra of \(\mathfrak{g}\) generated by \(\pmb{a}\) ; the elements of the kernel of \(f_{a}\) are called the relators of the family \(\pmb{a}\) . The family \(\pmb{a}\) is called generating (resp. free, basic) if \(f_{a}\) is surjective (resp. injective, bijective).

Let \(\mathfrak{g}\) be a Lie algebra. A presentation of \(\mathfrak{g}\) is an ordered pair \((\boldsymbol{a},\boldsymbol{r})\) consisting of a generating family \(\boldsymbol{a} = (a_{i})_{i\in \mathbb{I}}\) and a family \(\boldsymbol{r} = (r_j)_{j\in \mathbb{J}}\) of relators of \(\boldsymbol{a}\) generating the ideal of L(I) the kernel of \(f_{a}\) . We also say that \(\mathfrak{g}\) is presented by the family \(\boldsymbol{a}\) related by the relators \(r_j\) ( \(j\in J\) ).

Let I be a set and \(\pmb{r} = (r_j)_{j \in J}\) be a family of elements of the free Lie algebra L(I); let \(a_r\) be the ideal of L(I) generated by \(\pmb{r}\) . The quotient algebra L(I, \(\pmb{r}\) ) = L(I)/ \(a_r\) is called the Lie algebra defined by I and the family of relators ( \(r_j)_{j \in J}\) ; we also say that L(I, \(\pmb{r}\) ) is defined by the presentation (I, \(\pmb{r}\) ), or also by (I; ( \(r_j = 0\) ) \(_{j \in J}\) ). When the family \(\pmb{r}\) is empty, L(I, \(\pmb{r}\) ) = L(I).

Let \(I\) and \(r\) be as above; let \(\xi_i\) denote the image of \(i\) in \(L(I, r)\) . The generating family \(\xi = (\xi_i)_{i \in I}\) and the family of relators \(r\) constitute a presentation of \(L(I, r)\) . Conversely, if \(g\) is a Lie algebra and \((a, r)\) , where \(a = (a_i)_{i \in I}\) , is a presentation of \(g\) , there exists a unique isomorphism \(u: L(I, r) \to g\) such that \(u(\xi_i) = a_i\) for all \(i \in I\) .

\subsection{LIE POLYNOMIALS AND SUBSTITUTIONS}

Let I be a set. Let \(\mathrm{T}_i\) denote the canonical image of the element \(i\) of I in \(\mathbf{L}(\mathbf{I})\) (which is also sometimes denoted by \(\mathrm{L}((\mathrm{T}_i)_{i\in \mathbf{I}})\) ); the elements of \(\mathrm{L}(\mathbf{I})\) are called Lie polynomials in the indeterminates \((\mathrm{T}_i)_{i\in \mathbf{I}}\) .

Let g be a Lie algebra. If \(\boldsymbol{t} = (t_{i})_{i \in I}\) is a family of elements of g, let \(f_{\boldsymbol{t}}\) denote the homomorphism of L(I) into g such that \(f_{\boldsymbol{t}}(\mathrm{T}_{\boldsymbol{i}}) = t_{i}\) for \(i \in I\) (no. 2, Proposition 1). The image under \(f_{\boldsymbol{t}}\) of the element P of L(I) is denoted by \(\mathrm{P}((t_{i})_{i \in I})\) . In particular, \(\mathrm{P}((\mathrm{T}_{\boldsymbol{i}})_{i \in I}) = \mathrm{P}\) ; the above element \(\mathrm{P}((t_{i})_{i \in I})\) is sometimes called the element of g obtained by substituting the \(t_{i}\) for the \(T_{i}\) in the Lie polynomial \(\mathrm{P}((\mathrm{T}_{\boldsymbol{i}})_{i \in I})\) .

Let \(\sigma: \mathfrak{g} \to \mathfrak{g}'\) be a Lie algebra homomorphism. For every family \(t = (t_i)_{i \in I}\) of elements of \(\mathfrak{g}\) and all \(P \in L(I)\) ,

\[
\sigma (\mathrm{P} ((t _ {i}) _ {i \in \mathrm{I}})) = \mathrm{P} ((\sigma (t _ {i})) _ {i \in \mathrm{I}}),\tag{4}
\]

for \(\sigma^{\circ}f_{t}\) maps \(\mathrm{T}_i\) to \(\sigma(t_i)\) for \(i\in I\) .

Let \((\mathbf{Q}_j)_{j\in J}\) be a family of elements of \(\mathbf{L}(\mathbf{I})\) and let \(\mathbf{P}\in \mathbf{L}(\mathbf{J})\) . By substituting the \(Q_{j}\) for the \(T_{j}\) in \(P\) , we obtain a Lie polynomial \(R = P((Q_{j})_{j \in J}) \in L(I)\) . Then

\[
\mathrm{R} ((t _ {i}) _ {i \in \mathrm{I}}) = \mathrm{P} ((Q _ {j} ((t _ {i}) _ {i \in \mathrm{I}})) _ {j \in J})\tag{5}
\]

for every family \(\pmb{t} = (t_i)_{i\in \mathbf{I}}\) of elements of a Lie algebra \(\mathfrak{g}\) , as is seen by operating by the homomorphism \(f_{t}\) on the equation \(\mathrm{R} = \mathrm{P}((\mathrm{Q}_j)_{j\in \mathbf{J}})\) and using (4).

Let \(\mathfrak{g}\) be a Lie algebra, I a finite set and \(P \in L(I)\) . Suppose that \(\mathfrak{g}\) is a free K-module. The mapping

\[
\tilde {\mathrm{P}} \colon \mathfrak {g} ^ {\mathrm{I}} \to \mathfrak {g}
\]

defined by \(\tilde{\mathrm{P}}((t_i)_{i\in \mathbf{I}}) = \mathrm{P}((t_i)_{i\in \mathbf{I}})\) is then polynomial.† For the set F of mappings of \(\mathfrak{g}^{\mathbf{I}}\) into \(\mathfrak{g}\) is a Lie algebra with the bracket defined by

\[
[ \phi , \psi ] (\boldsymbol {t}) = [ \phi (\boldsymbol {t}), \psi (\boldsymbol {t}) ];\tag{6}
\]

the set \(\mathbf{F}'\) of polynomial mappings of \(\mathfrak{g}^{\intercal}\) into \(\mathfrak{g}\) is a Lie subalgebra of \(\mathbf{F}\) by the bilinearity of the bracket. Our assertion then follows from the fact that the mapping \(\mathrm{P} \mapsto \tilde{\mathrm{P}}\) is a Lie algebra homomorphism and \(\tilde{\mathrm{T}}_i = \mathrm{pr}_i \in \mathrm{F}'\) for all \(i\) .

\subsection{FUNCTORIAL PROPERTIES}

PROPOSITION 2. Let X and Y be two sets. Every mapping \(u: \mathrm{X} \to \mathrm{Y}\) can be extended uniquely to a Lie algebra homomorphism \(\mathrm{L}(u): \mathrm{L}(\mathrm{X}) \to \mathrm{L}(\mathrm{Y})\) . For every mapping \(v: \mathrm{Y} \to \mathrm{Z}\) , \(\mathrm{L}(v \circ u) = \mathrm{L}(v) \circ \mathrm{L}(u)\) .

The existence and uniqueness of \(\mathrm{L}(u)\) follow from Proposition 1 of no. 2. The homomorphisms \(\mathrm{L}(v\circ u)\) and \(\mathrm{L}(v)\circ \mathrm{L}(u)\) have the same restriction to X and hence are equal (Proposition 1).

COROLLARY. If \(u\) is injective (resp. surjective, bijective), so is \(L(u)\) .

Since the assertion is trivial for \(\mathrm{X} = \varnothing\) , we assume \(\mathrm{X} \neq \varnothing\) . If \(u\) is injective there exists a mapping \(v\) of \(\mathrm{Y}\) into \(\mathrm{X}\) such that \(v \circ u\) is the identity mapping of \(\mathrm{X}\) ; by Proposition 2, \(\mathrm{L}(v) \circ \mathrm{L}(u)\) is the identity of automorphism of \(\mathrm{L}(\mathrm{X})\) and hence \(\mathrm{L}(u)\) is injective. When \(u\) is surjective, there exists a mapping \(w\) of \(\mathrm{Y}\) into \(\mathrm{X}\) such that \(u \circ w\) is the identity mapping of \(\mathrm{Y}\) ; then \(\mathrm{L}(u) \circ \mathrm{L}(w)\) is the identity mapping of \(\mathrm{L}(\mathrm{Y})\) , which proves that \(\mathrm{L}(u)\) is surjective.

Let \(\mathbf{X}\) be a set and \(\mathbf{S}\) a subset of \(\mathbf{X}\) . The above corollary shows that the canonical injection of \(\mathbf{S}\) into \(\mathbf{X}\) can be extended to an isomorphism \(\alpha\) of \(\mathbf{L}(\mathbf{S})\) onto the Lie subalgebra \(\mathbf{L}'(\mathbf{S})\) of \(\mathbf{L}(\mathbf{X})\) generated by \(\mathbf{S}\) ; we shall identify \(\mathbf{L}(\mathbf{S})\) and \(\mathbf{L}'(\mathbf{S})\) by means of \(\alpha\) .

Let \((\mathrm{S}_{\alpha})_{\alpha \in \mathbf{I}}\) be a right directed family of subsets of X with union S. The relation \(\mathrm{S}_{\alpha} \subset \mathrm{S}_{\beta}\) implies \(\mathrm{L}(\mathrm{S}_{\alpha}) \subset \mathrm{L}(\mathrm{S}_{\beta})\) and hence the family of Lie subalgebras \(\mathrm{L}(\mathrm{S}_{\alpha})\) of \(\mathrm{L}(\mathrm{X})\) is right directed. Therefore \(\mathfrak{g} = \bigcup_{\alpha \in \mathbf{I}} \mathrm{L}(\mathrm{S}_{\alpha})\) is a Lie subalgebra of \(\mathrm{L}(\mathrm{X})\) ; then \(\mathrm{S} \subset \mathfrak{g}\) , whence \(\mathrm{L}(\mathrm{S}) \subset \mathfrak{g}\) , and, as \(\mathrm{L}(\mathrm{S}_{\alpha}) \subset \mathrm{L}(\mathrm{S})\) for all \(\alpha \in \mathbf{I}\) , \(\mathfrak{g} \subset \mathrm{L}(\mathrm{S})\) . Hence

\[
\mathrm{L} \left(\bigcup_ {\alpha \in \mathrm{I}} \mathrm{S} _ {\alpha}\right) = \bigcup_ {\alpha \in \mathrm{I}} \mathrm{L} (\mathrm{S} _ {\alpha})\tag{7}
\]

for every right directed family \((\mathrm{S}_{\alpha})_{\alpha\in\mathbb{I}}\) of subsets of X.

Applying the above to the family of finite subsets of X, we see that every element of \(\mathrm{L}(\mathrm{X})\) is of the form \(\mathrm{P}(x_{1},\ldots,x_{n})\) where P is a Lie polynomial in n indeterminates and \(x_{1},\ldots,x_{n}\) are elements of X.

PROPOSITION 3. Let \(\mathbf{K}'\) be a non-zero commutative ring and \(u: \mathbf{K} \to \mathbf{K}'\) a ring homomorphism. For every set \(\mathbf{X}\) there exists one and only one Lie \(\mathbf{K}'\) -algebra homomorphism

\[
v \colon \mathrm{L} _ {\mathbf {K}} (\mathrm{X}) \otimes \mathrm{K} ^ {\prime} \rightarrow \mathrm{L} _ {\mathbf {K} ^ {\prime}} (\mathrm{X})
\]

such that \(v(x \otimes 1) = x\) for \(x \in \mathbf{X}\) . Further, \(v\) is an isomorphism.

Applying Proposition 1 to \(\mathfrak{g} = \mathrm{L}_{\mathbb{K}'}(\mathrm{X})\) considered as a Lie K-algebra and the mapping \(x\mapsto x\) of \(\mathbf{X}\) into \(\mathfrak{g}\) , we obtain a K-homomorphism \(\mathrm{L}_{\mathbb{K}}(\mathrm{X})\to \mathrm{L}_{\mathbb{K}'}(\mathrm{X})\) , whence there is a K'-homomorphism \(v\colon \mathrm{L}_{\mathbb{K}}(\mathrm{X})\otimes \mathrm{K}'\to \mathrm{L}_{\mathbb{K}'}(\mathrm{X})\) . The fact that \(v\) is unique and is an isomorphism follows from the fact that the ordered pair \((\mathrm{L}_{\mathbb{K}}(\mathrm{X})\otimes \mathrm{K}',x\mapsto x\otimes 1)\) is a solution of the same universal mapping problem as the ordered pair \((\mathrm{L}_{\mathbb{K}'}(\mathrm{X}),x\mapsto x)\) .

Remark. Let \(\mathfrak{h}'\) be a Lie K'-algebra and \(\mathfrak{h}\) the Lie K-algebra derived from \(\mathfrak{h}'\) by restricting the ring of scalars. If \(\mathrm{P} \in \mathrm{L}_{\mathrm{K}}(\mathrm{X})\) , we can define \(\tilde{\mathrm{P}}: \mathfrak{h}^{\mathrm{X}} \to \mathfrak{h}\) (no. 4). We see immediately that

\[
\tilde {\mathrm{P}} = (v (\mathrm{P} \otimes 1)) ^ {\sim}.
\]
